\documentclass[final, 9pt, twocolumn]{saucis}
\usepackage{cite}
\usepackage{float}
\usepackage{epsfig}
\usepackage{epstopdf}
\usepackage{orcidlink}
\usepackage{rorlink}
\usepackage{url}
\let\CheckCommand\providecommand
\usepackage{microtype}
\usepackage{threeparttable}
\usepackage{xcolor}
\newcommand{\rev}[1]{\textcolor{black}{#1}}

\newcommand{\vo}{9} %volume
\newcommand{\n}{4} %number
\newcommand{\y}{2025} %year
\newcommand{\p}{345-432} %pages
\renewcommand{\saucisJD}{RESEARCH ARTICLE}
\newcommand{\dergiparkID}{19234532} %dergipark ID
\renewcommand{\saucisAuthors}{B. Koksoy, A. Pasaoglu, S. Duman and S. P. Iran} %author names
\startsAtpage{102} %page number start point

\begin{document}

%---------------------------------------------------------------------------------
%	TITLE
%---------------------------------------------------------------------------------

\title{Forecasting Provincial Healthcare Infrastructure and Budget Allocation in Turkey Using Machine Learning and TOPSIS}

%---------------------------------------------------------------------------------
%	AUTHORS
%---------------------------------------------------------------------------------
\author[1*]{Bedirhan Köksoy,\orcidlink{0009-0003-2358-9209}}
\author[2]{Ali Paşaoğlu,\orcidlink{0000-0002-6853-1356}}
\author[3]{Şevval Duman,\orcidlink{0009-0005-1777-4306}}
\author[4]{Sara Pournaji İran,\orcidlink{0009-0005-0029-6459}}
\affil[1]{Information Systems and Technologies
Department, Faculty of Computer and Information Sciences, Yeditepe University, Istanbul, Türkiye, \rorlink{https://ror.org/025mx2575}}
\affil[2]{Department of Computer Engineering, Faculty of Engineering and Natural Sciences, Istanbul Rumeli University, Istanbul, Türkiye,
\rorlink{https://ror.org/05av6y173}}
\affil[3]{Department of Electrical and Electronics Engineering, Faculty of Engineering, Istanbul University-Cerrahpasa, Istanbul, Türkiye,
\rorlink{https://ror.org/01dzn5f42}}
\affil[4]{Department of Management Information Systems, Faculty of Business Administration, Istanbul Rumeli University, Istanbul, Türkiye,
\rorlink{https://ror.org/05av6y173}}
\affil[*]{Bedirhan Köksoy}
\newcommand{\corresAuthoremail}{bedirhan.koksoy@yeditepe.edu.tr} %volume

%---------------------------------------------------------------------------------
%	ABSTRACT
%---------------------------------------------------------------------------------

\saucisAbstract{
Equitable allocation of healthcare resources across provinces requires accurate projections of future infrastructure needs, demographic pressures, and available budgets. This study proposes an integrated machine learning framework for province-level healthcare planning in Turkey covering the 2025--2030 period. The framework consists of three interconnected components. In the first component, six healthcare infrastructure indicators the number of "physicians, nurses, midwives, other healthcare personnel, hospital beds, and hospitals" are forecasted for all 81 Turkish provinces using 19 machine learning models trained on 2002--2024 data with a five-fold walk-forward cross-validation scheme. In the second component, four demographic and service capacity criteria "the 0--14 age group ratio, the 65+ age group ratio, the supply of physicians per 1,000 people, and the supply of hospital beds per 100,000 people" are projected to 2030 and used to construct TOPSIS-based provincial need rankings. Criterion weights are derived objectively using the Shannon entropy method. In the third component, Turkey's national health budget is forecasted in USD terms using an ensemble of time-series and regression models, and the results are converted to Turkish lira using projected exchange rates. The three components are integrated through a dynamic budget allocation mechanism with a feedback loop: after each annual allocation, provincial TOPSIS scores are updated to reflect received resources, ensuring that high-priority provinces in early years yield priority to underserved provinces in subsequent years. Results show that small northern provinces with high elderly population ratios and low physician density, such as Sinop and Artvin, rank highest in 2025, while the feedback mechanism redistributes priority to larger coastal provinces from 2026 onwards. The proposed framework provides a data-driven, adaptive, and reproducible approach to provincial health resource planning.
}
%---------------------------------------------------------------------------------
%	KEYWORDS 
%---------------------------------------------------------------------------------
\saucisKeywords{Budget Allocation, Healthcare Infrastructure Forecasting, Machine Learning, Shannon Entropy, TOPSIS}
%---------------------------------------------------------------------------------
%	ARTICLE INFO
%---------------------------------------------------------------------------------
\saucisreceived{23 November 2025}
\saucisaccepted{23 November 2025}
\saucisonline{23 November 2025}
\maketitle 

\section{Introduction}  %% Please avoid complex formulas in (sub)titles
Province-level healthcare resource planning in Turkey faces a structural challenge. Personnel assignments are still largely based on previous vacant position records, the departure of personnel who have completed their service period, or rotations initiated by provincial health directorates \cite{intro4}. Although this approach may address the personnel shortage of a particular hospital, it may also lead to workforce deficits in hospitals that have a greater need for personnel but are less preferred by healthcare workers. Budget management in healthcare is a fundamental component of health planning. One of the major challenges in Turkey's healthcare planning process is budget allocation. In general, a larger share of healthcare resources is allocated to metropolitan areas due to their higher population density. However, this situation may lead to financial constraints in smaller cities that also require adequate healthcare funding, thereby weakening residents' access to healthcare services \cite{intro5}. Turkey provides broad access to publicly financed healthcare \cite{intro3}, which places a substantial planning responsibility on the Ministry of Health, requiring forward-looking, data-driven tools rather than reactive workforce management.

The issue of healthcare planning in Turkey is addressed not only through the initiatives of the Ministry of Health but also through academic and independent studies. Various studies have been conducted on this topic, offering insights into how specific aspects of healthcare planning can be improved. Among these studies, some evaluate the distribution of healthcare personnel in Turkey, while others examine the economic dimension of healthcare planning \cite{intro6, intro7}.

Healthcare planning has gained a new perspective with the emergence of artificial intelligence (AI). In particular, AI-based forecasting and allocation approaches have become increasingly prominent in recent planning studies. Using AI-based methods, cities with greater personnel needs can be identified more effectively, and the criteria guiding budget allocation can be determined in a more systematic, data-driven manner.

In recent years, AI-based personnel allocation forecasting has become a prominent research topic in the literature. In one study, the physician supply and demand balance in Turkey for 2030 was modeled using XGBoost, yielding the most robust forecasts \cite{intro8}. In another, nurse staffing adequacy and patient care models in hospital wards were evaluated using machine learning, with Random Forest achieving the highest performance \cite{intro9}. However, both studies address workforce forecasting in isolation and do not link their projections to a need-based provincial ranking or a budget distribution mechanism.

In personnel forecasting, population balance is a key factor; however, demographic structure must also be evaluated independently of population size. The \rev{Technique for Order of Preference by Similarity to Ideal Solution (TOPSIS)} method offers a suitable framework for examining such regional and demographic differences. A TOPSIS-based study evaluating the sustainable development performance of TR33 region provinces across housing, transportation, urbanization, and environmental criteria demonstrated that multidimensional city-level ranking can be effectively applied in healthcare personnel planning as well \cite{intro10}.

\rev{Entropy-weighted TOPSIS has become an established approach in health
system evaluation. It has been used to assess the equity of regional health
resource allocation~\cite{lit3},~\cite{lit6}, to rank the service quality of
tertiary public hospitals at city level~\cite{lit4}, and to evaluate
community-level basic public health services~\cite{lit5}. At the cross-country
level, the same method has been applied to integrate equity measures across
demographic, geographic, and economic dimensions in ASEAN health
systems~\cite{lit2}, while in the Turkish context it was used to compare the
performance of 544 secondary-level public hospitals across geographical
regions~\cite{lit1}. These studies confirm that entropy-based objective
weighting is well suited to health-sector indicators, where subjective expert
weights are difficult to justify.}

\rev{Such evaluations, however, describe allocation as it currently stands
rather than projecting future need, and spatial inequalities in healthcare
resources have been shown to persist over time when allocation rules remain
fixed~\cite{lit7}.} Nevertheless, such ranking studies remain static and do not
incorporate forecasted infrastructure indicators or a dynamic resource feedback
mechanism across planning periods.

In healthcare planning, studies have used AI for budget planning. In one study, Turkey's per capita healthcare expenditures were estimated using AI techniques and genetic algorithm-based feature selection, yielding the best-performing results with the Random Forest model \cite{intro11}. In another study focusing on per capita healthcare expenditure, Turkey's healthcare expenditures were predicted using machine learning models, with the Support Vector Regression algorithm achieving the highest performance \cite{intro12}. Yet both studies operate at the national aggregate level and provide no mechanism for translating budget forecasts into province-level allocations.

This study extends our previous research, in which the regional distribution of healthcare workforce demand in Turkey was forecasted using ML algorithms based on 2024 personnel distribution patterns \cite{intro13}. However, the previous study did not incorporate demographic population characteristics, TOPSIS-based need prioritization, healthcare budget allocation, or improved machine learning methods. Therefore, the present study aims to provide a more comprehensive healthcare planning framework by integrating personnel distribution forecasts, demographic indicators, TOPSIS scores, and machine learning-based healthcare budget allocation for the 2025--2030 period. \rev{Table~\ref{tab:related_work} summarises the methodological scope of the studies discussed above and positions the present work with respect to them.}

\begin{table}[htbp]
\centering
\caption{\rev{Comparison of Related Studies and the Positioning of the Present Work}}
\label{tab:related_work}
\small
\begin{tabular}{p{1cm} p{3.3cm} p{3.0cm}}
\hline
\textbf{Study} & \textbf{Method and Scope} & \textbf{Missing Component} \\
\hline
\cite{intro8}  & XGBoost, national & No need ranking or allocation \\
\cite{intro9}  & Random Forest, hospital ward & No spatial dimension \\
\cite{intro11} & RF with genetic algorithm, national & Expenditure only \\
\cite{intro12} & SVR, GPR, DT, national & Expenditure only \\
\cite{intro10} & TOPSIS, regional & Static, observed data only \\
\cite{lit1}    & Entropy--TOPSIS, national & Static, retrospective \\
\cite{lit2}    & Entropy--TOPSIS with DEA, cross-country & Static, no forecasting \\
\cite{lit3}, \cite{lit6} & Entropy--TOPSIS, regional & Static, no budget link \\
\cite{intro13} & ML forecasting, province & No demographics, MCDM or budget \\
\hline
\textbf{This study} & 19 ML models, entropy--TOPSIS, dynamic feedback, province & --- \\
\hline
\end{tabular}
\end{table}

The main contributions of this study are as follows:

\begin{itemize}
    \item A province-level infrastructure forecast for all 81 Turkish provinces across six indicators using 19 machine learning models with walk-forward cross-validation.
    \item An entropy-weighted dynamic TOPSIS need ranking updated annually via a feedback loop that prevents budget monopolisation by early-priority provinces.
    \item A national health budget forecast in USD integrated with province-level allocations for the 2025--2030 period.
\end{itemize}

\section{Material and Methods}
In the relevant section of the study, the data sources used for the dataset, the preprocessing steps applied to the datasets, the pre-modeling preparation stages, the preferred machine learning (ML) approaches, weighted ensemble selection strategies, and the model evaluation metrics are explained in detail. The explanations are presented under separate headings to improve clarity.

\rev{The overall workflow of the proposed framework, together with the
data flow between its three modules, is presented in Figure~\ref{fig:workflow}.}

\begin{figure}[htbp]
    \centering
    \includegraphics[width=0.5\textwidth]{foto5.png}
    \caption{\rev{Overall workflow of the proposed healthcare planning framework,
    showing the infrastructure forecasting, need prioritization, and budget
    allocation modules and the dynamic feedback loop connecting them}}
    \label{fig:workflow}
\end{figure}

\subsection{Data Collection Sources}
This section describes the datasets used in the study. It explains the sources from which the datasets were collected and the types of data obtained from these sources. The subsequent analysis and processing steps applied to the collected data are described in detail in other sections.

\subsubsection{Dataset Sources}
The data used in this study were collected from three main source groups (Figure~\ref{fig1}). The primary source was the Turkish Statistical Institute (TURKSTAT), which provided demographic and healthcare-related statistics \cite{method1}. The second source was Investing.com, from which historical USD/TRY exchange rate data were obtained. The third source group was used to obtain actual and forecasted USD/TRY exchange rate data for future periods.

\begin{figure}[htbp]
		\centering
			\includegraphics[width=0.5\textwidth]{foto1.png} 
\caption{\rev{The Three Main Source Groups from Which the Datasets Were Collected}}
	\label{fig1}
\end{figure}

\subsubsection{Data Collected from the Data Sources}
After reliable data sources were identified in Section~2.1.1, the required datasets were collected in accordance with the study's intent. The collected data are summarized (Table~\ref{tab:sources1} and Table~\ref{tab:sources2}). The first table presents the source, spatial coverage, and year range of each dataset. In addition, the second table provides detailed information on the content and size of the corresponding data. The corresponding module associated with each dataset is also included in the tables. The purpose of each dataset and its role within the overall 
framework of the study are explained under Section 3 in the Results section.

\begin{table}[htbp]
\centering
\caption{Data Sources and Spatial Coverage of Collected Datasets}
\label{tab:sources1}
\small
\begin{tabular}{p{3cm} p{1.8cm} p{1.8cm}}
\hline
\textbf{Module} & \textbf{Data Source} & \textbf{Spatial Level} \\
\hline
Infrastructure Forecasting          & TURKSTAT      & Province \\
Need Prioritization    & TURKSTAT      & Province \\
Need Prioritization          & TURKSTAT      & Province \\
Need Prioritization          & TURKSTAT      & Province \\
Budget Allocation          & TURKSTAT      & Country  \\
Budget Allocation          & Investing.com & Country  \\
Budget Allocation          & Third Group   & Country  \\
\hline
\end{tabular}
\end{table}

Table~\ref{tab:sources1} shows that the datasets are drawn from two 
primary sources: TURKSTAT and external financial data providers. 
Province-level data cover the 81 provinces of Turkey and are used in the \rev{Infrastructure Forecasting Module (IFM)} and Need Prioritization modules, while country-level data support the macroeconomic analyses conducted in the Budget Allocation module. \rev{This division reflects the spatial resolution at which each quantity is defined, since infrastructure and demographic indicators vary across provinces while expenditure and exchange rate data are reported only at the national level.}

\begin{table}[htbp]
\centering
\caption{Year Range, Data Types, and Dataset Sizes}
\label{tab:sources2}
\small
\begin{tabular}{p{1.8cm} p{4.5cm} p{1.2cm}}
\hline
\textbf{Year Range} & \textbf{Data Types} & \textbf{Dataset Size} \\
\hline
2002--2024 & Population, number of hospitals, number of hospital beds, and healthcare personnel data & 13041 \\
2025, 2030 & Population data for 2025 and projected population data for 2030 & 162 \\
2007--2025 & Population data for the 0--14 and 65+ age groups & 3078 \\
2002--2024 & Supply of physicians per thousand people and supply of hospital beds per one hundred thousand people & 3726 \\
1999--2024 & General health expenditure data & 25 \\
1999--2024 & Monthly USD/TRY exchange rate data & 300 \\
2025--2030 & Actual and forecasted USD/TRY exchange rate data & 72 \\
\hline
\end{tabular}
\end{table}

Table~\ref{tab:sources2} presents the temporal scope and content of each dataset. \rev{The largest dataset comprises 13,041 data points, corresponding to 1,863 province-year observations across seven variables for the 2002--2024 period}, while the smallest comprises 25 annual country-level health expenditure records.

It should be noted that some datasets contain multiple sub-variables within a single data type. In particular, the healthcare personnel dataset comprises six distinct personnel categories: assistant physicians, specialist physicians, general practitioners, midwives, nurses, and other healthcare staff.

\subsection{Data Preprocessing}
Data preprocessing is the process of cleaning, transforming, and integrating raw data before analysis so that it can be processed more effectively and accurately by data mining algorithms. This process includes fundamental steps such as handling missing or noisy data through data cleaning, converting data into appropriate formats through data transformation, and removing unnecessary features from the dataset through data reduction. In this context, data preprocessing aims to maximize data quality and directly improve the accuracy, reliability, and performance of the ML or data mining models to be developed \cite{method2}.

In this study, the data preprocessing stage was carried out to enable machine learning algorithms to operate more effectively on the data. All procedures performed during the data preprocessing stage were organized into subsections according to the sources from which the data were obtained. Since no additional preprocessing was required for the Investing.com dataset group, preprocessing procedures were applied only to the TURKSTAT datasets and the third dataset group.

\subsubsection{TURKSTAT Dataset Preprocessing}

\begin{figure}[htbp]
		\centering
			\includegraphics[width=0.5\textwidth]{foto2.png} 
\caption{\rev{Merging of the Four TURKSTAT Datasets Used in the IFM}}
\label{fig2}
\end{figure}

TURKSTAT data were collected separately in Office Open XML Spreadsheet (XLSX) format from the web platform named Biruni, which provides access to TURKSTAT datasets \cite{method3}. The collected datasets were then converted into Comma-Separated Values (CSV) format, which can be processed more conveniently by machine learning models \cite{method4}. After the conversion process, the datasets were merged into a common file according to the study group in which they were used. For example, since \rev{the IFM used} data on the number of hospitals, hospital beds, healthcare personnel, and population, these four datasets were combined into a single file (Figure~\ref{fig2}).

\subsubsection{Third Group Dataset Preprocessing}
The third dataset group consisted of sources that provide USD/TRY exchange rate data as well as future exchange rate forecasts. The data in this group were directly collected in CSV format using a Python-based web scraper \cite{method5}. Subsequently, similar to the preprocessing procedure applied to the TURKSTAT datasets, the collected data were merged into a single CSV file in a suitable format to improve their usability for ML models \cite{method4}. Figure~\ref{fig3} illustrates the data extraction process using the web scraper and the integration of the collected data into a single CSV file.

\begin{figure}[htbp]
		\centering
			\includegraphics[width=0.5\textwidth]{foto3.png} 
\caption{\rev{Extraction and Aggregation of USD/TRY Exchange Rate Data Using a Python-Based Web Scraper}}
	\label{fig3}
\end{figure}

\subsection{Pre-Modeling Preparation}
Before proceeding to the prediction stage, several preparatory steps are required for ML models. These preparation stages are generally referred to as pre-modeling \cite{methodnew1}. In the pre-modeling stage, six different preparatory steps were performed. These steps are illustrated in Figure~\ref{fig4} and then explained in detail under separate subsections.

\begin{figure}[htbp]
		\centering
			\includegraphics[width=0.5\textwidth]{foto6.png} 
\caption{\rev{The Six Pre-Modeling Preparation Steps Applied Before Model Training}}
	\label{fig4}
\end{figure}

\subsubsection{Data Loading and Variable Derivation}

The dataset used in the IFM was compiled from province-level health records spanning the years 2002 to 2024, covering all 81 provinces of Turkey. The resulting dataset consists of $81 \times 23 = 1{,}863$ observations, which can be expressed as:

\begin{equation}
    N = P \times T = 81 \times 23 = 1{,}863
\end{equation}

where $N$ denotes the total number of observations, $P$ is the number of provinces, and $T$ is the number of years in the study period. In addition, a composite variable representing the total number of physicians was derived by aggregating three sub-categories of healthcare personnel available in the dataset. The derived variable is defined \rev{in Equation~2} as follows:

\begin{equation}
    \text{Physician} = \text{Specialist} + \text{General Practitioner} + \text{Resident}
\end{equation}

This aggregation was preferred over using each sub-category separately, as the distribution of residents across provinces exhibits considerable instability and abrupt fluctuations over time. The composite variable therefore provides a more stable and consistent representation of total physician availability at the province level.

\subsubsection{Feature Map}

Feature extraction is a critical process in the use of ML, 
as the selection of input variables directly influences model 
performance \cite{methodnew2}. In this study, a target-specific 
feature map was constructed, where each target variable is predicted using all remaining health infrastructure variables along with year and population as inputs. In other words, each target variable serves as an output while all other variables act as its predictors, ensuring that no target variable is used to predict itself. To illustrate the corresponding feature map structure applied in the study, the obtained feature assignments are presented in Figure ~\ref{fig5}.

\begin{figure}[htbp]
		\centering
			\includegraphics[width=0.5\textwidth]{foto7.png} 
\caption{\rev{Target-Specific Feature Map: Input Variables Assigned to Each Target}}
	\label{fig5}
\end{figure}

\subsubsection{Walk-Forward Cross-Validation}

Cross-validation is a resampling technique that estimates model 
generalization performance by partitioning data into $k$ subsets, 
using each as a test set once while training on the remainder 
\cite{methodnew3}. Empirical studies suggest that $k = 5$ or $k = 10$ generally yields reliable estimates \cite{methodnew4}.

However, standard random $k$-fold cross-validation is not suitable for time-series data, as randomly assigning observations to folds can cause future data to leak into the training set. To address this, a walk-forward strategy with an expanding training window was adopted, preserving the temporal order of observations across five folds as shown in Table~\ref{tab:kfold}.

\begin{table}[htbp]
\centering
\caption{\rev{Training and Test Periods of the Five Walk-Forward Folds}}
\label{tab:kfold}
\small
\begin{tabular}{p{1.5cm} p{3cm} p{3cm}}
\hline
\textbf{Fold} & \textbf{Training Period} & \textbf{Test Period} \\
\hline
1 & 2002--2016 & 2017--2020 \\
2 & 2002--2017 & 2018--2021 \\
3 & 2002--2018 & 2019--2022 \\
4 & 2002--2019 & 2020--2023 \\
5 & 2002--2020 & 2021--2024 \\
\hline
\end{tabular}
\end{table}

Fold~5 serves as the primary evaluation fold, as its test period 
covers the years closest to the forecast horizon of the study.

\subsubsection{Hyperparameter Optimization}

Hyperparameters are configuration settings that govern the learning 
process of a model and must be specified prior to training, as they cannot be learned directly from the data \cite{methodnew5}. The goal of hyperparameter optimization is to find the configuration $\lambda^{*}$ that minimizes a predefined loss function over a search space $\Lambda$:

\begin{equation}
    \lambda^{*} = \arg\min_{\lambda \in \Lambda} \; 
    \mathcal{L}\bigl(f_{\lambda}(X_{\text{train}}),\, X_{\text{val}}\bigr)
\end{equation}

where $f_{\lambda}$ denotes the model trained with hyperparameter 
configuration $\lambda$, and $\mathcal{L}$ is the loss function 
evaluated on the validation set.

In this study, Grid Search Cross-Validation (GridSearchCV) was 
employed, which exhaustively evaluates all possible combinations 
within a predefined hyperparameter grid \cite{methodnew6}. For a 
grid containing $m$ hyperparameters with $n_i$ candidate values 
for the $i$-th hyperparameter, the total number of configurations 
evaluated is:

\begin{equation}
    |\Lambda| = \prod_{i=1}^{m} n_i
\end{equation}

Each configuration $\lambda \in \Lambda$ is evaluated using 
$k$-fold cross-validation with $k = 2$ inner folds, and the 
mean validation loss across folds is computed as:

\begin{equation}
    \overline{\mathcal{L}}(\lambda) = 
    \frac{1}{k} \sum_{j=1}^{k} \mathcal{L}\bigl(f_{\lambda}^{(j)}\bigr)
\end{equation}

where $f_{\lambda}^{(j)}$ denotes the model trained on fold $j$. 
The scoring criterion used was the negative mean absolute error 
($-\text{MAE}$), so the optimal configuration satisfies:

\begin{equation}
    \lambda^{*} = \arg\min_{\lambda \in \Lambda} \; 
    \overline{\mathcal{L}}(\lambda) = 
    \arg\min_{\lambda \in \Lambda} \; 
    \frac{1}{k} \sum_{j=1}^{k} \frac{1}{n_j} 
    \sum_{i=1}^{n_j} \left| y_i - \hat{y}_i^{(\lambda, j)} \right|
\end{equation}

where $n_j$ is the number of observations in fold $j$, $y_i$ is 
the true value, and $\hat{y}_i^{(\lambda, j)}$ is the predicted 
value under configuration $\lambda$ in fold $j$.

\rev{The search grids and the resulting selections are summarised in
Table~\ref{tab:hyperparams}.}

\begin{table}[htbp]
\centering
\caption{\rev{Hyperparameter Search Grids and Values Selected at Fold~5}}
\label{tab:hyperparams}
\small
\begin{threeparttable}
\begin{tabular}{p{1cm} p{4cm} p{2.2cm}}
\hline
\textbf{Model} & \textbf{Search grid} & \textbf{Selected (Fold 5)} \\
\hline
LR & fit\_intercept: \{True, False\} & Target-specific \\
RF & n\_estimators: \{100, 200, 300\} \newline max\_depth: \{None, 10\} & 300 \newline 10 \\
ExtraTrees & n\_estimators: \{100, 200, 300\} \newline max\_depth: \{None, 10\} & Target-specific \newline 10 \\
GB & n\_estimators: \{100, 200\} \newline learning\_rate: \{0.05, 0.1\} \newline max\_depth: \{3, 5\} & 200 \newline 0.05 \newline 3 \\
XGBoost & n\_estimators: \{200, 300\} \newline learning\_rate: \{0.05, 0.1\} \newline max\_depth: \{3, 6\} & 200 \newline 0.05 \newline 3 \\
CatBoost & iterations: \{200, 300\} \newline depth: \{4, 6, 8\} & 300 \newline 8 \\
Ridge & Not tuned\tnote{a} & Default \\
ElasticNet & Not tuned\tnote{a} & Default \\
HGB & Not tuned\tnote{a} & Default \\
\hline
\end{tabular}
\begin{tablenotes}
\footnotesize
\item Selected values are the most frequently chosen setting across the six
target variables; ``target-specific'' indicates no single setting was dominant.
\item[a] Applied with the default settings of the corresponding library
implementation.
\end{tablenotes}
\end{threeparttable}
\end{table}
\rev{Grid search was applied with two inner folds, negative mean absolute error
as the scoring criterion, and a fixed random seed of 42 to ensure
reproducibility. Because optimisation was performed independently for each
target variable and each outer fold, the reported selections correspond to the
setting most frequently chosen across the six targets at Fold~5. Time-series and
baseline models were excluded from grid search, as their parameters are
estimated directly from the series rather than selected from a predefined grid.}

\subsubsection{Feature Scaling}

Feature scaling is a preprocessing step that brings all input 
variables to a common numerical range, which is particularly 
critical for regularized linear models such as Ridge, as the penalty terms of these models are sensitive to the magnitude input features \cite{methodnew7}. In this study, StandardScaler was applied exclusively to LR, and Ridge, while tree-based and boosting models were left unscaled as they are invariant to monotonic feature transformations. StandardScaler standardizes each feature by removing the mean and scaling to unit variance:

\begin{equation}
    z = \frac{x - \mu}{\sigma}
\end{equation}

where $x$ denotes original feature value, while $\mu$ and $\sigma$ represent the mean and standard deviation of training set. The scaler was fitted only on the training fold and then applied to test fold to avoid data leakage.

\subsubsection{Non-Negativity Constraint}

Since all target variables in this study represent counts of 
physical or human health resources, such as the number of 
hospitals, beds, or personnel, negative predicted values are 
inherently meaningless. To ensure the physical validity of all 
forecasts, a non-negativity constraint was applied by clipping 
all model outputs at zero, such that any predicted value 
$\hat{y} < 0$ was replaced by zero:

\begin{equation}
    \hat{y} = \max(0,\ \hat{y})
\end{equation}

This constraint was applied uniformly across all models and 
target variables following each prediction step.

\subsection{Machine Learning Approaches}
ML is the science of algorithms that enable computer systems to improve their task performance by learning from data-driven experience rather than being explicitly programmed. This discipline focuses on constructing statistical models capable of discovering structural patterns hidden within large datasets and making accurate predictions about future uncertainties or unseen cases. The main objective of this process is to enable a system to successfully analyze new examples that it has not previously encountered by applying the rules and generalizations learned from historical data \cite{method6}.

\begin{table}[htbp]
\centering
\caption{Machine Learning Model Groups and Their Assigned Modules}
\label{tab:models}
\small
\begin{threeparttable}
\begin{tabular}{p{1cm} p{2.5cm} p{4.2cm}}
\hline
\textbf{Module} & \textbf{Main Category} & \textbf{Models Included} \\
\hline
M1, M2       & Classical Regression Models  & LinearRegression (LR), Ridge, ElasticNet \\
M1, M2, M3   & Tree-Based Models            & RandomForest (RF), ExtraTrees, GradientBoosting (GB), HistGradientBoosting (HGB), XGBoost, CatBoost \\
M1, M2, M3   & Time-Series Models           & ARIMA, Prophet, HoltWinters \\
M3           & Time-Series Models           & Holt, HoltDamped \\
M1, M2       & Ensemble Models              & Voting, Weighted Ensemble \\
M2, M3       & Baseline/Trend Models        & Baseline\_Lag1, CAGR5, CAGR10 \\
\hline
\end{tabular}
\begin{tablenotes}
\footnotesize
\item \rev{M1: Infrastructure Forecasting Module (IFM), M2: Need Prioritization, M3: Budget Allocation. These labels and the module names are used interchangeably.}
\end{tablenotes}
\end{threeparttable}
\end{table}

In this study, a total of nineteen different ML algorithms were used across three modules. These algorithms were employed to compare the prediction results and thereby identify the approach that produced the highest prediction performance. The ML algorithms used in the study are presented in Table~\ref{tab:models}
according to their names and the corresponding module.

\subsubsection{Classical Regression Models}

Classical regression methods are statistical techniques used to model the linear association between a response variable and one or more explanatory variables by estimating a parametric function from the available data. These models assume that the relationship between variables can be adequately described by a linear combination of predictors and rely on well-established estimation techniques. Despite their relative simplicity, they offer high interpretability and serve as important benchmarks against which more complex methods can be evaluated. Three classical regression models are employed in this study.

\begin{itemize}
    \item Linear Regression estimates the dependent variable as a 
    linear combination of independent variables using ordinary least 
    squares (OLS), which minimizes the sum of squared residuals 
    \cite{method7}.
    \item Ridge Regression extends OLS by adding an L2 penalty term to reduce multicollinearity. It shrinks coefficients toward zero without setting them exactly to zero \cite{method8}.
    \item ElasticNet combines L1 and L2 penalty terms, offering both variable selection and multicollinearity control, merging the advantages of Ridge and LASSO \cite{method8}.
\end{itemize}

The advantages and disadvantages of the classical regression models are presented in Table~\ref{tab:classical_pros_cons} .

\begin{table}[htbp]
\centering
\caption{Advantages and Disadvantages of Classical Regression Models}
\label{tab:classical_pros_cons}
\small
\begin{tabular}{p{1.2cm} p{3.15cm} p{3.15cm}}
\hline
Model & Advantages & Disadvantages \\
\hline
Linear Regression
& Simple to interpret; low computational cost
& Sensitive to multicollinearity and outliers \\
Ridge
& Controls multicollinearity; retains all variables
& No variable selection; coefficients never reach zero \\
ElasticNet
& Performs variable selection and controls multicollinearity
& Requires tuning of two hyperparameters \\
\hline
\end{tabular}
\end{table}

\subsubsection{Tree-Based Models}

Tree-based models are non-parametric ensemble methods that partition the feature space through a series of recursive binary splits. Each split is determined by selecting the feature and threshold that best separates the target variable, allowing the model to approximate complex, non-linear decision boundaries without requiring feature scaling or distributional assumptions. These models have been shown to consistently outperform deep learning approaches on typical tabular datasets, owing to their robustness to irrelevant features, their ability to handle mixed data types, and their capacity to capture interaction effects without extensive preprocessing \cite{method10}.

\begin{itemize}
    \item RF constructs multiple decision trees using bootstrap sampling and random feature subsets, then combines their predictions to reduce variance and improve generalization ability. \cite{method10}.

    \item ExtraTrees extends the random forest idea by 
    selecting split thresholds randomly rather than optimally, which 
    further reduces variance at the cost of a slight increase in bias 
    \cite{method11}.

    \item GB generates trees sequentially; each tree corrects the remaining errors of the previous tree by following the negative slope of the loss function.
    \cite{method10}.

    \item HGB is a histogram-based variant 
    of GB that discretizes continuous features into bins 
    before splitting, substantially reducing training time on large 
    datasets \cite{method10}.

    \item XGBoost is a regularized and computationally 
    optimized GB framework that supports parallel 
    processing and natively handles missing values \cite{method10}.

    \item CatBoost is a GB algorithm 
    specifically designed to handle categorical features without manual encoding, using ordered boosting to prevent target leakage 
    \cite{method12}.
\end{itemize}

The advantages and disadvantages of the tree-based models are presented in Table~\ref{tab:tree_pros_cons} .

\begin{table}[htbp]
\centering
\caption{Advantages and Disadvantages of Tree-Based Models}
\label{tab:tree_pros_cons}
\small
\begin{tabular}{p{1.2cm} p{3.15cm} p{3.15cm}}
\hline
Model & Advantages & Disadvantages \\
\hline
RF
& Low variance; robust to noise and outliers
& High memory usage; slow on large datasets \\
ExtraTrees
& Faster training than RF; further variance reduction
& Slightly higher bias; less interpretable \\
GB & High predictive accuracy; handles mixed data types
& Prone to overfitting; slow sequential training \\
HGB
& Fast on large datasets; native missing value support
& Less interpretable; requires bin size tuning \\
XGBoost
& Regularization reduces overfitting; parallelizable
& Many hyperparameters; sensitive to tuning \\
CatBoost
& Native categorical support; robust to overfitting
& Slower training than XGBoost; high memory usage \\
\hline
\end{tabular}
\end{table}

\subsubsection{Time-Series Models}

Time-series models are developed to represent temporal patterns and trend structures in ordered sequential data. Unlike cross-sectional methods, these models explicitly account for the order of observations and the patterns that evolve over time. Five time-series models are employed in this study.

\begin{itemize}
    \item ARIMA (Autoregressive Integrated Moving Average) 
    models a time series through autoregressive and moving average 
    components after differencing to achieve stationarity 
    \cite{method13}.

    \item Prophet is a decomposable forecasting model 
    developed by Meta that separately models trend, seasonality, 
    and holiday effects, making it robust to missing data and 
    structural breaks \cite{method13}.

    \item Holt-Winters extends simple exponential smoothing 
    by incorporating both trend and seasonal components through 
    double or triple smoothing equations \cite{method14}.

    \item Holt applies double exponential smoothing to 
    capture level and trend in non-seasonal time series, using 
    two smoothing parameters \cite{method14}.

    \item Holt Damped introduces a damping parameter to 
    the Holt method, moderating the trend component over the 
    forecast horizon to avoid over-optimistic long-range predictions 
    \cite{method14}.
\end{itemize}

\begin{table}[htbp]
\centering
\caption{Advantages and Disadvantages of Time-Series Models}
\label{tab:timeseries_pros_cons}
\small
\begin{tabular}{p{0.9cm} p{3.3cm} p{3.3cm}}
\hline
Model & Advantages & Disadvantages \\
\hline
ARIMA
& Flexible for stationary series; well-established statistical foundation
& Requires stationarity; struggles with non-linear patterns \\
Prophet
& Handles missing data and outliers; easy to configure
& Less effective on short or irregular series \\
Holt-Winters
& Captures both trend and seasonality; interpretable parameters
& Assumes fixed seasonal pattern; sensitive to outliers \\
Holt
& Simple and effective for trended non-seasonal data
& Cannot capture seasonality; may overfit short series \\
Holt Damped
& Reduces over-forecasting in long horizons; more conservative
& Requires tuning of the damping parameter \\
\hline
\end{tabular}
\end{table}

The advantages and disadvantages of the time-series models are presented in Table~\ref{tab:timeseries_pros_cons} .

\subsubsection{Ensemble Models}

Ensemble models integrate the outputs of several base learners to improve generalization compared with an individual model. By combining models with different characteristics, ensemble techniques can lower variance, bias, or both \cite{method15}. In this study, two ensemble approaches are used.

\begin{itemize}
    \item Voting Ensemble aggregates the predictions of 
    multiple independently trained models by averaging their outputs. In this study, RF, XGBoost, and GB are combined using soft voting \cite{method15}.

    \item Weighted Ensemble assigns optimized weights 
    to each base model's prediction rather than treating all 
    models equally. Weights are determined by minimizing 
    validation error, allowing better-performing models to 
    contribute more to the final prediction \cite{method16}.
\end{itemize}

The advantages and disadvantages of the ensemble models are presented in Table~\ref{tab:ensemble_pros_cons} .

\begin{table}[htbp]
\centering
\caption{Advantages and Disadvantages of Ensemble Models}
\label{tab:ensemble_pros_cons}
\small
\begin{tabular}{p{1.2cm} p{3.15cm} p{3.15cm}}
\hline
Model & Advantages & Disadvantages \\
\hline
Voting
& Simple to implement; reduces individual model errors
& Equal weighting may dilute strong models \\
Weighted Ensemble
& Optimized contribution of each model; adaptive
& Weight optimization requires careful validation strategy \\
\hline
\end{tabular}
\end{table}

\subsubsection{Baseline and Trend Models}

Baseline and trend models provide simple yet informative reference points against which the performance of more complex models can be assessed. Rather than learning intricate patterns from data, these models rely on straightforward statistical rules or growth rate calculations. They are particularly useful for evaluating whether sophisticated algorithms offer meaningful improvements over naive forecasting strategies. Three baseline and trend models are employed in this study.

\begin{itemize}
    \item Baseline Lag1 is a naive forecasting model that 
    uses the most recent observed value as the prediction for the 
    next period. Despite its simplicity, it serves as a fundamental benchmark in time-series forecasting tasks \cite{method17}.

    \item \textbf{CAGR5} estimates future values by applying the 
    compound annual growth rate calculated over the preceding five 
    years. This approach captures medium-term growth momentum and 
    provides a trend-based projection without model fitting 
    \cite{method18}.

    \item \textbf{CAGR10} follows the same logic as CAGR5 but 
    extends the calculation window to ten years, offering a 
    longer-term perspective on growth trends and reducing the 
    influence of short-term fluctuations \cite{method18}.
\end{itemize}

\begin{table}[htbp]
\centering
\caption{Advantages and Disadvantages of Baseline and Trend Models}
\label{tab:baseline_pros_cons}
\small
\begin{tabular}{p{1.0cm} p{3.25cm} p{3.25cm}}
\hline
Model & Advantages & Disadvantages \\
\hline
Baseline Lag1
& Simple to implement; requires no training
& Ignores long-term trends; poor on volatile series \\
CAGR5
& Captures medium-term growth; interpretable
& Sensitive to the chosen base year; assumes constant growth \\
CAGR10
& Smooths short-term fluctuations; stable long-term estimate
& May miss recent structural changes; assumes constant growth \\
\hline
\end{tabular}
\end{table}

The advantages and disadvantages of the baseline and trend models are presented in Table~\ref{tab:baseline_pros_cons} .

\subsection{Weighted Ensemble Selection Strategy}

While Section~2.4.4 introduced the general ensemble modeling 
approaches employed in this study, the present section describes 
the specific ensemble selection strategy applied during the 
model evaluation phase. Rather than using a fixed ensemble 
configuration, a data-driven weighted ensemble was constructed 
separately for each target variable based on cross-validation 
performance.

Ensemble methods combine the predictions of multiple base learners to achieve lower generalization error than any individual model alone \cite{methodnew8}. The underlying principle is that diverse models tend to make different errors, and their aggregation reduces the overall prediction variance \cite{methodnew9}. Ensemble approaches have been shown to consistently outperform single models across a wide range of regression and classification tasks \cite{methodnew10}.

In this study, for each target variable $y_k$, the top-$K$ models 
were selected based on their Fold-5 accuracy scores, and their 
predictions were combined through a normalized weighted average. 
Let $\mathcal{M}_k = \{m_1, m_2, \ldots, m_K\}$ denote the set of 
top-$K$ models selected for target $y_k$. The weight assigned to 
model $m_i$ \rev{is given in Equation~9}:

\begin{equation}
    w_i = \frac{\text{Accuracy}_{i}}{\displaystyle\sum_{j=1}^{K} 
    \text{Accuracy}_{j}}, \quad \sum_{i=1}^{K} w_i = 1
\end{equation}

The ensemble prediction for target $y_k$ is then computed as:

\begin{equation}
    \hat{y}_k^{\text{ens}} = \sum_{i=1}^{K} w_i \cdot \hat{y}_k^{(m_i)}
\end{equation}

where $\hat{y}_k^{(m_i)}$ is the prediction of model $m_i$ for 
target $y_k$, and $K = 3$ in this study.

\subsection{Model Evaluation Metrics}
In this study, various metrics were used to compare the performance of the machine learning models. Through these metrics, the best-performing prediction model was selected, enabling future forecasts to be generated using the most reliable model. The evaluation metrics used in the study can be categorized into two main groups. The first group consists of basic regression metrics, while the second group consists of additional classification-based metrics. The contents of these metric groups are visualized in Figure~\ref{fig6} and explained in detail under the following subsections.

\begin{figure}[htbp]
		\centering
			\includegraphics[width=0.5\textwidth]{foto4.png} 
\caption{\rev{Regression and Classification-Based Metrics Used for Model Evaluation}}
	\label{fig6}
\end{figure}

In the \rev{Results} section of the article, all metric scores, including Accuracy, MAE, RMSE, R², Error Rate, Precision, Recall, and F1-score, are reported. However, during the visualization stage, only the most important metrics were included in the tables to avoid excessive complexity. The complete source code and detailed results are available for overall evaluation.

\subsubsection{Basic Regression Metrics}
The performance of all regression models was assessed using five 
widely adopted error metrics. Following \cite{method19}, the 
coefficient of determination ($R^{2}$) is considered the most 
informative among standard regression evaluation criteria, as it 
provides a normalized, scale-independent measure of explanatory power.

\begin{itemize}
    \item MAE measures the average 
    magnitude of prediction errors without considering their direction. It is expressed in the same unit as the target variable and is robust to large individual errors \cite{method19}.
    \begin{equation}
        \text{MAE} = \frac{1}{n} \sum_{i=1}^{n} \left| y_i - \hat{y}_i \right|
    \end{equation}

    \item Root Mean Square Error (RMSE) penalizes larger errors more severely than MAE due to the squaring process, making it more sensitive to outliers \cite{method19}.
    \begin{equation}
        \text{RMSE} = \sqrt{\frac{1}{n} \sum_{i=1}^{n} \left( y_i - \hat{y}_i \right)^{2}}
    \end{equation}

    \item Mean Absolute Percentage Error (MAPE) expresses 
    the prediction error as a percentage of the actual values, enabling 
    scale-independent comparisons across different targets \cite{method19}.
    \begin{equation}
        \text{MAPE} = \frac{1}{n} \sum_{i=1}^{n} 
        \left| \frac{y_i - \hat{y}_i}{\max\left(|y_i|,\, \varepsilon\right)} \right| \times 100
    \end{equation}
    where $\varepsilon$ is a small constant introduced to prevent 
    division by zero.

    \item Coefficient of Determination ($R^{2}$) quantifies 
    the proportion of variance in the target variable explained by the model. A value of 1 indicates a perfect fit, while values close to 0 or negative indicate poor predictive performance 
    \cite{method19}.
    \begin{equation}
        R^{2} = 1 - \frac{\sum_{i=1}^{n}(y_i - \hat{y}_i)^{2}}
        {\sum_{i=1}^{n}(y_i - \bar{y})^{2}}
    \end{equation}
    where $\bar{y}$ denotes the mean of observed values.

    \item Accuracy is derived from MAPE and represents the 
    percentage of prediction correctness. It is bounded below 
    by zero to avoid negative accuracy values \cite{method19}.
    \begin{equation}
        \text{Accuracy} = \max(0,\ 100 - \text{MAPE})
    \end{equation}
\end{itemize}

The properties and interpretation of each evaluation metric are 
summarized in Table~\ref{tab:metrics_summary}.

\begin{table}[htbp]
\centering
\caption{Summary of Regression Evaluation Metrics}
\label{tab:metrics_summary}
\small
\begin{tabular}{p{0.9cm} p{1cm} p{1cm} p{4.3cm}}
\hline
Metric & Range & Optimal Value & Interpretation \\
\hline
MAE      & $[0, +\infty)$ & $\rightarrow 0$   & Average absolute error; same unit as target \\
RMSE     & $[0, +\infty)$ & $\rightarrow 0$   & Penalizes large errors; sensitive to outliers \\
MAPE     & $[0, +\infty)$ & $\rightarrow 0$   & Scale-free percentage error \\
$R^{2}$  & $(-\infty, 1]$ & $\rightarrow 1$   & Proportion of variance explained by the model \\
Accuracy & $[0, 100]$     & $\rightarrow 100$ & Percentage correctness derived from MAPE \\
\hline
\end{tabular}
\end{table}

\subsubsection{Classification-Based Additional Metrics}

In addition to regression metrics, classification-based metrics were computed by binarizing the regression outputs to assess the directional accuracy of the forecasts. Specifically, predictions and actual values were converted into binary classes based on whether the value increased or decreased relative to the previous period. Precision, Recall, and F1-score were then applied to evaluate the model's ability to correctly identify the direction of change \cite{method20}.

\begin{itemize}
    \item Precision indicates the share of true positive predictions among all samples classified as positive. A higher precision score means that the model makes fewer false positive predictions \cite{method20}.
    \begin{equation}
        \text{Precision} = \frac{TP}{TP + FP}
    \end{equation}

    \item Recall represents the fraction of actual positive instances that are correctly recognized by the model. It shows how effectively the model captures true positives while minimizing missed positive cases \cite{method20}.
    \begin{equation}
        \text{Recall} = \frac{TP}{TP + FN}
    \end{equation}

    \item F1-Score is the harmonic mean of Precision and 
    Recall, providing a single balanced measure that accounts for 
    both false positives and false negatives \cite{method20}.
    \begin{equation}
        \text{F1} = 2 \times \frac{\text{Precision} \times \text{Recall}}
        {\text{Precision} + \text{Recall}}
    \end{equation}
\end{itemize}

where $TP$, $FP$, and $FN$ denote true positives, false positives, and false negatives, respectively. The summary of the classification-based evaluation metrics is presented in Table~\ref{tab:cls_metrics_summary}.

\begin{table}[htbp]
\centering
\caption{Summary of Classification-Based Evaluation Metrics}
\label{tab:cls_metrics_summary}
\small
\begin{tabular}{p{0.9cm} p{1.0cm} p{1cm} p{4.3cm}}
\hline
Metric & Range & Optimal Value & Interpretation \\
\hline
Precision & $[0, 1]$ & $\rightarrow 1$ & Fraction of correct positive predictions \\
Recall    & $[0, 1]$ & $\rightarrow 1$ & Fraction of actual positives correctly detected \\
F1-Score  & $[0, 1]$ & $\rightarrow 1$ & Harmonic mean of Precision and Recall \\
\hline
\end{tabular}
\end{table}

\section{Results}
In this section of the article, the procedures carried out throughout the study
are described in detail. The presentation is divided into three main parts, which
represent the core stages of the study. \rev{The overall workflow of these three
parts was presented in Figure~\ref{fig:workflow} in Section~2.} The outputs
obtained from the model predictions are evaluated in the Discussion section,
where the findings are assessed within the overall framework of the study.

\subsection{Health Infrastructure Forecast}

In the IFM of the study, healthcare-related data for the period between 2025 and 2030 were forecasted using machine learning algorithms based on health data from 2002 to 2024. \rev{The module proceeds in five stages: the province-level records are assembled and preprocessed, candidate models are trained and compared on the 2002--2024 period, the best performers for each target are combined into a weighted ensemble, population values for the intermediate forecast years are filled by extrapolation, and the ensemble is finally applied to produce the 2025--2030 projections.} The procedures carried out in this part are explained in detail under the following subsections. The flowchart of the IFM is presented in Figure~\ref{fig9}.
\begin{figure}[htbp]
		\centering
			\includegraphics[width=0.5\textwidth]{foto9.png} 
\caption{\rev{Five-Stage Workflow of the Infrastructure Forecasting Module}}
\label{fig9}
\end{figure}

\subsubsection{Data Preparation}

In the first stage of this part, the province-level data on population, number of hospitals, number of hospital beds, and healthcare personnel, as described in the “Data Collection Sources” section, were collected for the period between 2002 and 2024. \rev{Since the dataset, comprising 13,041 data points across 1,863 province-year observations, was stored in separate files}, these files were merged into a single CSV file as explained in the ``Data Preprocessing'' section.

\subsubsection{Model Development}

In the second stage of this part, the data from 2002 to 2024 were used to develop the best-performing prediction model. The purpose of developing these models was to identify the machine learning algorithm that could most accurately predict the available data and, accordingly, generate reliable forecasts for the period between 2025 and 2030.

Before proceeding to the modeling stage, the preparation steps described under the “Pre-Modeling Preparation” section were applied to the data, making the dataset ready for model development. Subsequently, the machine learning methods specified under the “Machine Learning Approaches” section and identified as relevant to the first part of the study were employed. The performance results of the models are presented in Table~\ref{tab:model_comparison}.

\begin{table}[htbp]
\begin{threeparttable}
\centering
\caption{\rev{Comparison of Machine Learning Models for Infrastructure Forecasting, Averaged Across Six Target Variables and Five Walk-Forward Folds}}
\label{tab:model_comparison}
\small
\begin{tabular}{p{1.8cm} p{1.2cm} p{1.2cm} p{1.3cm} p{1.1cm}}
\hline
Model & Accuracy (\%) & MAE & RMSE & $R^{2}$ \\
\hline
\textbf{HoltWinters} & \textbf{91.594} & \textbf{176.692} & \textbf{573.802} & \textbf{0.972} \\
ARIMA                & 90.886 & 221.379 &  768.090 & 0.947 \\
Prophet              & 90.527 & 202.769 &  684.702 & 0.963 \\
Voting               & 86.059 & 247.361 &  820.054 & 0.941 \\
ExtraTrees           & 85.708 & 240.752 &  759.110 & 0.951 \\
CatBoost             & 85.553 & 247.939 &  823.359 & 0.935 \\
XGBoost              & 85.447 & 264.961 &  895.543 & 0.922 \\
HGB & 85.269 & 335.039 & 1426.567 & 0.829 \\
RF         & 85.185 & 263.475 &  847.860 & 0.938 \\
GB     & 85.073 & 250.023 &  757.372 & 0.948 \\
LR     & 77.584 & 254.076 &  526.481 & 0.940 \\
Ridge                & 76.197 & 256.111 &  526.497 & 0.940 \\
ElasticNet           & 75.101 & 277.180 &  601.941 & 0.937 \\
\hline
\end{tabular}
\begin{tablenotes}
\footnotesize
\item \rev{Best-performing values are shown in bold.}
\end{tablenotes}
\end{threeparttable}
\end{table}

\subsubsection{Ensemble Model}

The results obtained under Section 3.1.2 represent the best overall performance averaged across all target variables. However, this averaging introduces a limitation: a model that predicts the physician variable very accurately but performs poorly on the hospital variable will have a reduced average score and may not be ranked as the best model overall, even though it is highly effective for a specific target. To address this, the three best-performing models were selected separately for each target variable and combined into a weighted ensemble using the approach described in Section 2.5. The resulting ensemble weights and model selections for each target are presented in Table~\ref{tab:ensemble_selection}.

\begin{table}[htbp]
\begin{threeparttable}
\centering
\caption{\rev{Three Best-Performing Models Selected for Each Target Variable and Their Normalised Ensemble Weights (Fold 5)}}
\label{tab:ensemble_selection}
\small
\begin{tabular}{p{0.8cm} p{1.2cm} p{1cm} p{1.2cm} p{1.2cm} p{0.8cm}}
\hline
\textbf{Target} & \textbf{Model} & \textbf{Accuracy} & \textbf{MAE} & \textbf{RMSE} & \textbf{Weight} \\
\hline
\textbf{Other}    & \textbf{ARIMA}                & \textbf{90.598} & \textbf{375.373} & \textbf{1146.703} & \textbf{0.337} \\
Other    & HoltWinters          & 90.479 & 379.220 & 1206.848 & 0.336 \\
Other    & HGB & 87.873 & 607.133 & 2742.473 & 0.327 \\
\hline
\textbf{Physician}   & \textbf{HoltWinters}          & \textbf{91.976} & \textbf{231.136} &  \textbf{685.106} & \textbf{0.337} \\
Physician   & ARIMA                & 90.619 & 300.441 &  947.616 & 0.332 \\
Physician   & CatBoost             & 90.058 & 283.047 & 1082.126 & 0.330 \\
\hline
\textbf{Midwife}  & \textbf{ARIMA}                & \textbf{93.777} &  \textbf{34.045} &   \textbf{68.273} & \textbf{0.346} \\
Midwife  & HoltWinters          & 92.300 &  39.515 &   64.822 & 0.340 \\
Midwife  & RF         & 85.265 &  82.051 &  127.816 & 0.314 \\
\hline
\textbf{Hospital} & \textbf{ARIMA}                & \textbf{95.956} &   \textbf{0.654} &    \textbf{1.411} & \textbf{0.356} \\
Hospital & HoltWinters          & 93.197 &   0.975 &    1.572 & 0.346 \\
Hospital & Voting               & 80.074 &   2.558 &    5.838 & 0.297 \\
\hline
\textbf{Nurse}    & \textbf{ExtraTrees}           & \textbf{93.937} & \textbf{166.521} &  \textbf{421.617} & \textbf{0.335} \\
Nurse    & Voting               & 93.531 & 198.314 &  551.330 & 0.333 \\
Nurse    & XGBoost              & 93.184 & 219.763 &  624.091 & 0.332 \\
\hline
\textbf{Bed}      & \textbf{ARIMA}                & \textbf{92.710} & \textbf{264.546} &  \textbf{837.160} & \textbf{0.341} \\
Bed      & HoltWinters          & 92.033 & 263.947 &  882.778 & 0.339 \\
Bed      & CatBoost             & 86.801 & 273.435 &  530.161 & 0.320 \\
\hline
\end{tabular}
\begin{tablenotes}
\footnotesize
\item \rev{Best-performing values are shown in bold.}
\end{tablenotes}
\end{threeparttable}
\end{table}

Following Table~\ref{tab:ensemble_selection}, the performance metrics of the developed community model are given in Table ~\ref{tab:ensemble_metrics}.

\begin{table}[htbp]
\begin{threeparttable}
\centering
\caption{\rev{Performance of the Weighted Ensemble by Target Variable}}
\label{tab:ensemble_metrics}
\small
\begin{tabular}{p{1.2cm} p{1.2cm} p{1.2cm} p{1.2cm} p{1.2cm}}
\hline
\textbf{Target} & \textbf{Accuracy} & \textbf{MAE} & \textbf{RMSE} & \textbf{$R^{2}$} \\
\hline
Other    & 90.814 & 354.194 & 1003.492 & 0.966 \\
Physician   & 92.499 & 258.361 &  875.994 & 0.977 \\
Midwife  & 91.851 &  42.386 &   66.102 & 0.995 \\
\textbf{Hospital} & \textbf{92.527} &   \textbf{0.949} &    \textbf{1.580} & \textbf{0.997} \\
\textbf{Nurse}    & \textbf{93.910} & \textbf{183.656} &  \textbf{512.585} & \textbf{0.991} \\
Bed      & 92.722 & 207.703 &  531.969 & 0.992 \\
\hline
\end{tabular}
\begin{tablenotes}
\footnotesize
\item \rev{Best-performing values are shown in bold.}
\end{tablenotes}
\end{threeparttable}
\end{table}

It was observed that ensemble modeling produced lower results than individual models for some predictions. However, despite this, it was decided to proceed with the ensemble model. The reason for this decision was the risk of flat-forecast behaviour.

\subsubsection{Flat Forecast Behaviour of ARIMA on Sparse Series}
When ARIMA encounters a near-constant time series, it collapses 
to a random walk model and produces flat forecasts equal to the 
last observed value \cite{result1}. In such cases, the model 
may achieve an artificially high accuracy on historical test data, while failing to generate meaningful forecasts beyond the training horizon \cite{result2}.

For example, if no new hospital was built in Bolu between 2021 and 2024, and ARIMA correctly predicted the hospital number in 2021, it could continue to predict the same unchanged value until 2024 and therefore achieve an apparently perfect prediction score (Table~\ref{tab:arima_trap}). However, when future forecasts are generated for the 2025–2030 period, ARIMA may continue to produce constant values, which can lead to inaccurate long-term forecasts. For this reason, ensemble modeling was constructed to generate more dynamic and variable predictions.

\begin{table}[htbp]
\centering
\caption{\rev{Example of Flat ARIMA Forecasts: Hospital Count in Bolu Province, 2021--2024}}
\label{tab:arima_trap}
\small
\begin{tabular}{p{1cm} p{1cm} p{1cm} p{1cm} p{1cm}}
\hline
\textbf{Province} & \textbf{Year} & \textbf{Target} & \textbf{Actual} & \textbf{Predicted} \\
\hline
Bolu & 2021 & Hospital & 11.0 & 11.0 \\
Bolu & 2022 & Hospital & 11.0 & 11.0 \\
Bolu & 2023 & Hospital & 11.0 & 11.0 \\
Bolu & 2024 & Hospital & 11.0 & 11.0 \\
\hline
\end{tabular}
\end{table}

\subsubsection{Population Projections for 2025--2030}

After developing an ensemble model that predicted data for the period 2002--2024 with a mean accuracy of 92\%, the next step was to forecast health infrastructure variables for 2025--2030. However, this required province-level population values for each year in the forecast horizon as input features.

As stated in Table~\ref{tab:sources1} under Section~2.1.2, TURKSTAT provides province-level population projections only for the reference years 2025 and 2030. Since intermediate years were not directly available, population values for 2026--2029 were estimated using linear extrapolation based on the historical trend of the last ten years of observed data. For each province $p$, a linear trend model \rev{shown in Equation~19} was fitted using ordinary least squares over the most recent ten observations:

\begin{equation}
    \hat{P}_{p}(t) = a_p \cdot t + b_p
\end{equation}

where $t$ denotes the year, and the coefficients $a_p$ and $b_p$ are estimated by minimizing the sum of squared residuals over the training window:

\begin{equation}
    [a_p,\, b_p] = \arg\min_{a,b} \sum_{i=1}^{10} 
    \Bigl(P_{p,i} - (a \cdot t_i + b)\Bigr)^{2}
\end{equation}

The fitted model was then evaluated at each target year $t \in \{2026, 2027, 2028, 2029\}$ to produce the projected population values. A total of 486 province-year pairs were filled using this extrapolation procedure.

\subsubsection{Health Infrastructure Forecast for 2025--2030}

Using the previously established weighted ensemble model and the 
province-level population projections for 2025--2030, health 
infrastructure variables were forecast at the province level for 
each year in the target horizon. The forecast was produced for 
six target variables: number of hospitals, hospital beds, physicians, nurses, midwives, and other healthcare personnel.

The forecasting procedure was applied sequentially for each year 
$t \in \{2025, 2026, \ldots, 2030\}$. Since each target variable 
uses the remaining targets as input features, a circular dependency arises during future-year prediction. To resolve this, an iterative update strategy was adopted: for each forecast year, the prediction loop was executed $r = 3$ times, allowing the estimates to converge as each target's predicted value is progressively refined. Formally, at iteration $r$ and year $t$ \rev{as given in Equation~21}:

\begin{equation}
    \hat{y}_k^{(r)}(t) = \sum_{i=1}^{3} w_i \cdot 
    f_i\!\left(\hat{y}_{-k}^{(r-1)}(t),\, P(t)\right)
\end{equation}

where $\hat{y}_{-k}^{(r-1)}(t)$ denotes the predicted values of 
all targets except $y_k$ from the previous iteration, $P(t)$ is 
the projected population for year $t$, and $f_i$ is the $i$-th 
model in the ensemble. All predictions were clipped at zero and 
rounded to the nearest integer. The final output consists of 
$81 \times 6 = 486$ province-year observations covering the 
2025--2030 period.

\subsection{TOPSIS-Based Provincial Ranking}

\rev{TOPSIS ranks alternatives according to their geometric closeness to a positive ideal solution (PIS) and separation from a negative ideal solution (NIS)}~\cite{result3}. The best alternative is the one located nearest to the PIS while remaining farthest from the NIS. The method has been widely applied in healthcare facility planning and resource allocation problems \cite{result4}.

\begin{table}[h]
\centering
\caption{\rev{The Four Criteria Used in the TOPSIS Ranking and Their Optimisation Direction}}
\label{tab:topsis_criteria}
\begin{threeparttable}
\begin{tabular}{p{3cm} p{1cm} p{2.2cm} p{1cm}}
\hline
\textbf{Criterion} & \textbf{Unit} & \textbf{\rev{Interpretation}} & \textbf{Type} \\
\hline
0--14 age group population ratio & \%          & Higher $\rightarrow$ more demand    & Benefit \\
65+ age group population ratio   & \%          & Higher $\rightarrow$ more demand    & Benefit \\
\rev{Physicians} per 1,000 people & per 1,000   & Higher $\rightarrow$ better supply  & Cost \\
Hospital beds per 100,000 people & per 100,000 & Higher $\rightarrow$ better supply  & Cost \\
\hline
\end{tabular}
\begin{tablenotes}
\footnotesize
\item \rev{Supply indicators are treated as cost criteria because higher supply implies lower relative need.}
\end{tablenotes}
\end{threeparttable}
\end{table}

Provincial rankings were constructed using four criteria that reflect both demographic pressure and service capacity, as shown in Table~\ref{tab:topsis_criteria}. Criterion weights were determined objectively using the Shannon entropy method, which identifies the most discriminating criterion based on the distribution of values across provinces \cite{result3}. The TOPSIS procedure was applied as follows:

\begin{enumerate}
    \item Data collection and preprocessing
    \item Entropy-based criterion weighting
    \item Model development using the available data
    \item Future forecasting using the developed model
    \item Sequential Resource Allocation with Feedback
\end{enumerate}

\subsubsection{Data Collection and Preprocessing}

During the data collection stage, the 0--14 and 65+ population data for the period between 2007 and 2025 were obtained from TURKSTAT, as described in Section~2.1. In addition, the number of physicians per 1,000 people and the number of hospital beds per 100,000 people for the period between 2002 and 2024 were collected. Province-level population data for the period between 2002 and 2024 were also included. Following the data collection process, the preprocessing steps described in Section~2.2 were applied, and the files were combined into a single CSV file. The number of data records used for the TOPSIS model is presented in Table~\ref{tab:topsis_data}.

\begin{table}[h]
\centering
\caption{Number of Data Records Used for the TOPSIS Model}
\label{tab:topsis_data}
\begin{tabular}{p{4.5cm} p{1.5cm} p{1cm}}
\hline
\textbf{Variable} & \textbf{Period} & \textbf{Records} \\
\hline
0--14 age group population ratio    & 2007--2025 & 1{,}539 \\
65+ age group population ratio      & 2007--2025 & 1{,}539 \\
Physicians per 1,000 people            & 2002--2024 & 1{,}863 \\
Hospital beds per 100,000 people    & 2002--2024 & 1{,}863 \\
Province-level total population     & 2002--2024 & 1{,}863 \\
\hline
\textbf{Total}                      &            & \textbf{8{,}667} \\
\hline
\end{tabular}
\end{table}

\subsubsection{Entropy-Based Criterion Weighting}

Shannon entropy, originally introduced in the context of information theory, quantifies the uncertainty or information content within a probability distribution \cite{result5, result6}. A criterion with low entropy exhibits high variability across alternatives, meaning it carries more discriminating power. Conversely, a criterion with high entropy shows little variation and contributes less to the ranking \cite{result7}. This property makes Shannon entropy a widely used objective weighting method in multi-criteria decision-making, as it derives weights directly from the data without requiring subjective judgment.

In this study, criterion weights were calculated from the 2025 provincial data using the following procedure. Let $x_{ij}$ denote the value of criterion $j$ for province $i$, and let $m$ be the number of provinces ($m = 81$).

Step 1 — Normalization \rev{(Equation~22)}:
\begin{equation}
p_{ij} = \frac{x_{ij}}{\sum_{i=1}^{m} x_{ij}}
\end{equation}

Step 2 — Entropy calculation:
\begin{equation}
E_j = -k \sum_{i=1}^{m} p_{ij} \ln(p_{ij}), \quad k = \frac{1}{\ln(m)}
\end{equation}

Step 3 — Divergence:
\begin{equation}
d_j = 1 - E_j
\end{equation}

Step 4 — Weight assignment:
\begin{equation}
w_j = \frac{d_j}{\sum_{j=1}^{n} d_j}
\end{equation}

The resulting weights are presented in Table~\ref{tab:topsis_weights}. The 65+ age ratio received the highest weight (37.0\%), reflecting the greatest inter-provincial variability and thus the strongest discriminating power among the four criteria.

\begin{table}[h]
\centering
\caption{Entropy-Based Criterion Weights}
\label{tab:topsis_weights}
\begin{tabular}{p{4.5cm} p{1.5cm} p{1cm}}
\hline
\textbf{Criterion} & \textbf{Weight} & \textbf{(\%)} \\
\hline
65+ age group population ratio   & 0.3698 & 37.0 \\
0--14 age group population ratio & 0.1976 & 19.8 \\
Physicians per 1,000 people         & 0.2205 & 22.1 \\
Hospital beds per 100,000 people & 0.2121 & 21.2 \\
\hline
\textbf{Total}                   & \textbf{1.0000} & \textbf{100.0} \\
\hline
\end{tabular}
\end{table}

\subsubsection{Model Development Using the Available Data}

The same ML models described in Section~2.4 were applied to the four TOPSIS criteria using a five-fold walk-forward cross-validation scheme. The primary objective was to identify the best-performing model for each criterion and subsequently generate forecasts through 2030. Each target variable was trained on its own available historical range: 2007--2025 for age group ratios and 2002--2024 for the clinical supply indicators.

\rev{Initial results were strong for two criteria, with Ridge selected for the
65+ age ratio and LR for physicians per 1,000. The 0--14 age ratio and the
hospital bed ratio, however, yielded low or negative $R^{2}$ values across all
models. For these two targets the feature set was therefore revised: lagged
values of the target were added and any variable used to construct the target
was excluded to prevent leakage. Baseline\_Lag1, a naive predictor carrying
forward the previous year's value, was also added to the candidate pool as a
reference. Because of this revision, the two targets were evaluated on the three
most recent years (2022--2024) rather than the five-fold walk-forward scheme
used for the other criteria. HGB then performed best for the 0--14 age ratio and
ElasticNet for the hospital bed ratio, both exceeding the Baseline\_Lag1
reference. Table~\ref{tab:topsis_model_best} presents the final model selected
for each criterion.}

\begin{table}[h]
\centering
\caption{\rev{Best-Performing Model Selected for Each TOPSIS Criterion
(Fold~5 for the first two criteria; 2022--2024 evaluation for the last two)}}
\label{tab:topsis_model_best}
\begin{tabular}{p{2cm} p{1.1cm} p{0.8cm} p{0.8cm} p{0.8cm} p{0.8cm} p{0.8cm}}
\hline
\textbf{Target} & \textbf{Model} & \textbf{Acc. (\%)} & \textbf{MAE} & \textbf{RMSE} & \textbf{R\textsuperscript{2}} \\
\hline
65+ age ratio        & Ridge                & 99.28 & 0.072  & 0.085  & 0.864 \\
Physicians per 1,000    & LR     & 98.94 & 0.026  & 0.027  & 0.954 \\
0--14 age ratio      & HGB & 99.39 & 0.128  & 0.191  & 0.999 \\
Beds per 100,000     & ElasticNet           & 98.70 & 3.540  & 4.280  & 0.997 \\
\hline
\end{tabular}
\end{table}

\subsubsection{Future Forecasting Using the Developed Models}
Using the best-performing model identified for each criterion, provincial forecasts were generated for the 2025--2030 period. Ridge was used for the 65+ age ratio, LR for physicians per 1,000, HGB for the 0--14 age ratio, and ElasticNet for hospital beds per 100,000.

\subsubsection{Sequential Resource Allocation with Feedback}
A critical issue arises when TOPSIS scores are computed statically across the entire 2025--2030 horizon: a province with persistently high need indicators would receive the largest budget share every year, regardless of whether prior allocations have already begun to address its deficiencies. To overcome this, a dynamic feedback mechanism was adopted. After each annual allocation, the TOPSIS criteria values are updated to reflect resources already directed to each province, and the ranking is recalculated for the following year. This ensures that provinces whose needs have been partially met yield priority to others with greater remaining deficiencies, as illustrated in Figure~\ref{fig10}.

\begin{figure}[htbp]
		\centering
			\includegraphics[width=0.5\textwidth]{foto10.png} 
\caption{\rev{Static Versus Dynamic TOPSIS Allocation: Fixed Ranking Compared with Annual Recalculation}}
	\label{fig10}
\end{figure}

\subsection{Health Budget Modeling}

The third component of this study focuses on forecasting Turkey's national health budget for the 2025--2030 period. The forecasts were produced in USD to avoid distortions caused by Turkish lira depreciation, and subsequently converted to TRY using projected exchange rates. The subsections below detail the data collection, exchange rate construction, model development, and future budget forecasting steps. The workflow of the Health Budget Forecast section is shown in Figure~\ref{fig11}.

\begin{figure}[htbp]
		\centering
			\includegraphics[width=0.5\textwidth]{foto11.png} 
\caption{\rev{Four-Stage Workflow of the Health Budget Forecasting Module}}
	\label{fig11}
\end{figure}

\subsubsection{Data Collection and Preprocessing}

In this section, the data collection and preprocessing steps for budget forecasting are described. As outlined in Section~2.1, the dataset consists of Turkey's annual general health expenditure figures for the period 1999--2024 and the monthly USD-TRY exchange rate for the same period. Following data collection, the preprocessing steps described in Section~2.2 were applied. The number of records used is presented in Table~\ref{tab:budget_data}.

\begin{table}[h]
\centering
\caption{Number of Data Records Used for Budget Modeling}
\label{tab:budget_data}
\begin{tabular}{p{4.5cm} p{1.5cm} p{1cm}}
\hline
\textbf{Variable} & \textbf{Period} & \textbf{Records} \\
\hline
Turkey annual health expenditure & 1999--2024 & 25 \\
Monthly USD/TRY exchange rate    & 1999--2024 & 300 \\
\hline
\textbf{Total} & & \textbf{325} \\
\hline
\end{tabular}
\end{table}

\subsubsection{USD/TRY Exchange Rate Construction for 2025--2030}

Since forecasting Turkey's health budget directly in Turkish lira would introduce distortions due to currency depreciation, all budget figures were converted to USD, requiring the construction of a monthly USD/TRY exchange rate series for the 2025--2030 period. Data collection process is described in Section~2.1, and Table~\ref{tab:usdtry_sources} summarizes the number of actual versus estimated monthly observations by year.

\begin{table}[h]
\centering
\caption{Monthly USD/TRY Data Composition by Year}
\label{tab:usdtry_sources}
\begin{tabular}{p{1cm} p{1cm} p{1cm} p{4cm}}
\hline
\textbf{Year} & \textbf{Actual (months)} & \textbf{Forecast (months)} & \textbf{Source} \\
\hline
2025 & 12 & 0  & X-Rates, ExchangeRates.org \\
2026 & 5  & 7  & Mixed \\
2027 & 0  & 12 & LongForecast + interpolation \\
2028 & 0  & 12 & LongForecast + interpolation \\
2029 & 0  & 12 & LongForecast + interpolation \\
2030 & 0  & 12 & LongForecast + interpolation \\
\hline
\textbf{Total} & \textbf{17} & \textbf{55} & \\
\hline
\end{tabular}
\end{table}

The resulting annual average USD/TRY rates used in budget conversion are presented in Table~\ref{tab:usdtry_annual}.

\begin{table}[h]
\centering
\caption{Annual Average USD/TRY Exchange Rates (2025--2030)}
\label{tab:usdtry_annual}
\begin{tabular}{lc}
\hline
\textbf{Year} & \textbf{USD/TRY} \\
\hline
2025 & 39.54 \\
2026 & 46.32 \\
2027 & 55.93 \\
2028 & 66.08 \\
2029 & 83.17 \\
2030 & 96.18 \\
\hline
\end{tabular}
\end{table}

\subsubsection{Model Development Using 1999--2024 Data}

The health expenditure figures were extracted from the 1999--2024 dataset, comprising annual total health expenditure, per capita expenditure, and the public and private sector expenditure ratios. Since the Turkish lira has undergone significant depreciation over this period, forecasting directly in TRY would inflate apparent growth rates. Therefore, all budget figures were converted to USD using the corresponding annual average exchange rate. Several time-series features were derived to preserve the temporal structure of the data, including the previous year's budget, the two-year lagged budget, a rolling moving average, year-on-year growth rates, and a linear trend index. Eight models were trained on this dataset, covering time-series methods (ARIMA, Holt, HoltDamped), compound annual growth rate models (CAGR5, CAGR10), and tree-based regressors (RF, ExtraTrees, GB). Model performance was evaluated using a within-sample test scheme, with particular emphasis on forecast accuracy in the most recent years \rev{to assess forward-looking reliability, as reported in Table~\ref{tab:budget_model_general}}.

\begin{table}[h]
\begin{threeparttable}
\centering
\caption{\rev{Comparison of Models for National Health Budget Forecasting (USD Target)}}
\label{tab:budget_model_general}
\begin{tabular}{lccc}
\hline
\textbf{Model} & \textbf{Accuracy (\%)} & \textbf{MAE} & \textbf{RMSE} \\
\hline
\textbf{ExtraTrees}       & \textbf{90.210} & \textbf{4{,}965}  & \textbf{6{,}157}  \\
CAGR10           & 88.370 & 5{,}571  & 6{,}826  \\
GB & 88.254 & 5{,}662  & 6{,}884  \\
CAGR5            & 88.153 & 5{,}633  & 6{,}962  \\
ARIMA            & 88.114 & 5{,}559  & 6{,}618  \\
HoltDamped          & 86.677 & 6{,}011  & 7{,}148  \\
RF     & 86.786 & 6{,}127  & 7{,}349  \\
Holt    & 86.119 & 6{,}256  & 7{,}407  \\
\hline
\end{tabular}
\begin{tablenotes}
\footnotesize
\item \rev{Best-performing values are shown in bold.}
\end{tablenotes}
\end{threeparttable}
\end{table}

Rather than relying on a single best-performing model, the top five models were combined into a weighted ensemble to reduce the risk of overfitting to any individual model's strengths and to improve forecast stability. Ensemble weights were assigned proportionally to each model's Fold-5 accuracy, so that higher-performing models contributed more to the final prediction. The resulting ensemble included a diverse mix of tree-based regressors, a time-series model, and growth rate estimators, providing complementary perspectives on the budget trajectory. The selected models and their ensemble weights are shown in Table~\ref{tab:budget_ensemble}.

\begin{table}[h]
\begin{threeparttable}
\centering
\caption{\rev{Models Selected for the Budget Forecasting Ensemble and Their Weights}}
\label{tab:budget_ensemble}
\begin{tabular}{lcc}
\hline
\textbf{Model} & \textbf{Accuracy (\%)} & \textbf{Ensemble Weight} \\
\hline
\textbf{ExtraTrees}       & \textbf{90.210} & \textbf{0.2201} \\
CAGR10           & 88.370 & 0.1962 \\
GB & 88.254 & 0.1930 \\
CAGR5            & 88.153 & 0.1940 \\
ARIMA            & 88.114 & 0.1966 \\
\hline
\end{tabular}
\begin{tablenotes}
\footnotesize
\item \rev{Best-performing values are shown in bold.}
\end{tablenotes}
\end{threeparttable}
\end{table}

\subsubsection{Budget Forecasting for 2025--2030}

The weighted ensemble model constructed in the previous section was applied to generate Turkey's national health budget forecasts for the 2025--2030 period. Since forecasting directly in Turkish lira would introduce distortions due to currency depreciation, all projections were produced in USD. The ensemble combined the outputs of ExtraTrees, ARIMA, CAGR10, CAGR5, and GB models, with weights proportional to each model's Fold-5 accuracy. The resulting USD figures were subsequently converted to TRY using the projected annual average exchange rates presented in Table~\ref{tab:usdtry_annual}.

\subsection{Integration of TOPSIS Rankings and Budget Allocation}
In the final stage, the 2025--2030 TOPSIS provincial priority rankings were combined with the forecasted national health budget. Rather than applying a static allocation based on a fixed ranking, a dynamic feedback mechanism was adopted. After each annual budget distribution, the received funds are treated as a partial improvement in a province's healthcare capacity, thereby reducing its need score in subsequent years. To quantify this effect, a budget factor with a multiplier of 0.25 was introduced, meaning that prior allocations account for 25\% of the updated TOPSIS score in each following year. The value of 0.25 was selected to reflect a moderate adjustment rate, ensuring that prior allocations meaningfully reduce a province's relative need score without dominating the demographic and supply indicators that form the core of the TOPSIS ranking. This corresponds to a weighting scheme in which received resources account for one quarter of the updated score, preserving the primacy of structural need over historical funding patterns. This ensures that provinces which have already received substantial funding yield priority to others with greater remaining deficiencies. The complete allocation cycle is illustrated in Figure~\ref{fig12}.

\begin{figure}[h]
    \centering
		\includegraphics[width=0.5\textwidth]{foto12.png} 
\caption{\rev{Annual Cycle of the Dynamic Budget Allocation Mechanism}}
\label{fig12}
\end{figure}

\subsection{Evaluation of Results}

The dynamic TOPSIS-based allocation yielded different ranks for the provinces for the period 2025–2030, which indicates that the feedback mechanism is capable of avoiding a fixed set of provinces to dominate the budgetary resources year after year.

In 2025, the prioritized provinces were mostly small cities in the north of Anatolia such as Sinop, Artvin, Kastamonu, Çankırı and Bartın. The common profile of these provinces is that they have high ratios of 65 years old and above and low physician density, which is why they record high TOPSIS scores. Their small populations also resulted in very high per capita allocations (Artvin and Sinop approximately 282,800 TRY and 214,000 TRY per person respectively in 2025).

These large per-capita transfers had the direct effect of providing the top-ranked provinces with enough resources in 2025 to reduce their relative need scores. From 2026 on, the ranking changed significantly and Balıkesir rose to the first position, where it stayed until 2030, with Aydın, Muğla, Manisa and Mersin in the upper ranks. These are larger coastal provinces with aging populations and moderate clinical supply indicators that were previously displaced by smaller, more severely underserved cities. Their larger populations result in lower per capita allocations.

Table~\ref{tab:topsis_top5} summarises the top five provinces by TOPSIS priority rank for each year of the forecast horizon. The dynamic feedback mechanism causes a notable structural shift after the first forecast year: while small northern provinces dominate in 2025 due to high elderly population ratios and low supply coverage, accumulated budget support from 2026 onward redistributes relative need, bringing larger western coastal provinces to the top of the priority list. Provinces receiving large early investments thus progressively yield their top positions to others with greater remaining deficiencies, which aligns with the core objective of equitable and need-responsive resource distribution.

\begin{table}[ht]
\centering
\caption{Top 5 Priority Provinces by Year, Dynamic TOPSIS (2025--2030)}
\label{tab:topsis_top5}
\small
\begin{threeparttable}
\begin{tabular}{p{0.6cm} p{1.5cm} p{1.7cm} p{2.2cm}}
\hline
\textbf{Year} & \textbf{Province} & \textbf{TOPSIS Score} & \textbf{Per Capita (TRY)} \\
\hline
\multirow{5}{*}{2025}
 & \textbf{Sinop} & \textbf{0.712} & \textbf{214{,}167} \\
 & Artvin    & 0.697 & 282{,}800 \\
 & Kastamonu & 0.689 & 123{,}239 \\
 & Çankırı   & 0.689 & 233{,}376 \\
 & Bartın    & 0.661 & 217{,}190 \\
\hline
\multirow{5}{*}{2026}
 & \textbf{Balıkesir} & \textbf{0.826} & \textbf{40{,}089} \\
 & Aydın     & 0.810 & 42{,}911 \\
 & Ordu      & 0.806 & 65{,}728 \\
 & Muğla     & 0.799 & 44{,}663 \\
 & Manisa    & 0.797 & 33{,}646 \\
\hline
\multirow{5}{*}{2027}
 & \textbf{Balıkesir} & \textbf{0.794} & \textbf{51{,}987} \\
 & Aydın     & 0.774 & 55{,}248 \\
 & Ordu      & 0.763 & 84{,}487 \\
 & Muğla     & 0.761 & 56{,}739 \\
 & Manisa    & 0.760 & 43{,}397 \\
\hline
\multirow{5}{*}{2028}
 & \textbf{Balıkesir} & \textbf{0.782} & \textbf{64{,}016} \\
 & Aydın     & 0.762 & 67{,}848 \\
 & Muğla     & 0.750 & 69{,}017 \\
 & Manisa    & 0.748 & 53{,}507 \\
 & Ordu      & 0.747 & 104{,}047 \\
\hline
\multirow{5}{*}{2029}
 & \textbf{Balıkesir} & \textbf{0.780} & \textbf{83{,}548} \\
 & Aydın     & 0.759 & 88{,}248 \\
 & Muğla     & 0.748 & 88{,}925 \\
 & Manisa    & 0.745 & 69{,}917 \\
 & Mersin    & 0.742 & 51{,}187 \\
\hline
\multirow{5}{*}{2030}
 & \textbf{Balıkesir} & \textbf{0.781} & \textbf{100{,}262} \\
 & Aydın     & 0.758 & 105{,}496 \\
 & Muğla     & 0.749 & 105{,}401 \\
 & Manisa    & 0.745 & 84{,}049 \\
 & Mersin    & 0.744 & 61{,}288 \\
\hline
\end{tabular}
\begin{tablenotes}
\footnotesize
\item \rev{Within each year, provinces are listed in descending order of TOPSIS score, so the first row corresponds to the province with the greatest assessed need. The highest-ranked province of each year is shown in bold.}
\end{tablenotes}
\end{threeparttable}
\end{table}

\section{\rev{Discussion}}
\subsection{\rev{Comparison with the Literature}}

\rev{The forecasting results of this study diverge from those reported in
comparable work. Previous studies on healthcare workforce and expenditure
forecasting in Turkey identified tree-based or kernel-based learners as the
strongest performers: XGBoost for physician supply and demand~\cite{intro8},
Random Forest for nurse staffing~\cite{intro9} and for per capita health
expenditure~\cite{intro11}, and Support Vector Regression for national health
spending~\cite{intro12}. Here, time-series methods outperformed these families
across all six infrastructure targets, with HoltWinters, ARIMA, and Prophet
reaching 91.59\%, 90.89\%, and 90.53\% accuracy respectively, against 85.71\%
for the best tree-based model (Table~\ref{tab:model_comparison}).}

\rev{This divergence follows from the structure of the data. Each province-level
series contains only 23 annual observations and exhibits a strong monotonic
trend driven by sustained capacity expansion since 2002. Tree-based learners
predict piecewise-constant values within observed ranges and cannot extrapolate
beyond the training range, a limitation visible in the lowest coefficient of
determination in the comparison ($R^{2} = 0.829$ for HGB). The cited studies
operate on cross-sectional data or on national aggregates over longer horizons,
conditions under which this constraint does not bind, so the present findings
complement rather than contradict them. Time-series methods carry their own
failure mode on near-constant series, which the weighted ensemble was adopted to
contain (Section~3.1.4).}

\rev{A second body of work applies entropy-weighted TOPSIS to health system
evaluation~\cite{lit1,lit2,lit3,lit4,lit5,lit6}. These studies establish that
entropy-based objective weighting suits health-sector indicators, and the
present study adopts the same weighting principle. They differ from it in
temporal orientation: each ranks units on observed data at a single point in
time, and the resulting ranking does not respond to any allocation that follows
from it, as is also the case for the provincial ranking in~\cite{intro10}. The
consequence has been documented directly, since spatial inequalities in hospital
bed distribution across 2308 Chinese counties persist where allocation rules do
not adjust to prior investment~\cite{lit7}. The feedback mechanism proposed here
addresses this condition by reducing a province's need score in proportion to
funds already received. Table~\ref{tab:static_dynamic} quantifies its effect.}

\begin{table}[htbp]
\centering
\caption{\rev{Effect of the Feedback Mechanism on the Circulation of Provincial Priority (2025--2030)}}
\label{tab:static_dynamic}
\small
\begin{tabular}{p{3.6cm} p{1.4cm} p{1.6cm}}
\hline
\textbf{Indicator} & \textbf{Static} & \textbf{Dynamic} \\
\hline
Distinct provinces entering the top five & 5 & 11 \\
Years before the top five first changes & --- & 1 \\
Provinces ranked first at any point & 1 & 2 \\
\hline
\end{tabular}
\end{table}

\rev{Under a static ranking computed once on 2025 data, the same five provinces
would hold the top five positions in every year of the horizon. Under the
dynamic scheme, eleven distinct provinces enter the top five and the composition
changes structurally after the first allocation year, as the small northern
provinces prioritised in 2025 yield to larger western coastal provinces whose
deficiencies remain unaddressed. Persistence is reduced rather than eliminated:
Balıkesir remains first from 2026 to 2030, reflecting a demographic and supply
profile that a single feedback cycle does not resolve.}

\subsection{\rev{Practical Implications}}

\rev{The entropy weighting assigned 37.0\% of the total criterion weight to the
65+ population ratio, against 19.8\% for the 0--14 ratio and a combined 43.3\%
for the two supply indicators (Table~\ref{tab:topsis_weights}). Because entropy
weights derive from inter-provincial variability rather than expert judgement,
this indicates that the elderly population share is currently the strongest
discriminator among Turkish provinces. For planning purposes it implies that
allocation criteria weighted primarily by total population or by absolute
personnel deficits will systematically understate need in provinces undergoing
demographic ageing, a pattern concentrated in the northern and central Anatolian
provinces that dominate the 2025 ranking. The per capita figures reported in
Table~\ref{tab:topsis_top5}  express relative
priority rather than operationally feasible transfers, and in practice would
require bounding by absorptive capacity; the framework accommodates such
constraints without structural modification, since the allocation stage is
separable from the ranking stage.}

\subsection{Limitations and Future Work}
The proposed framework yields promising results in the context of province level healthcare resource planning, although some limitations need to be taken into account. First, the exchange rate projections used to express the budget forecasts in Turkish lira were obtained from third party forecasting sources for the 2026--2030 period. Since long-term currency forecasts involve considerable uncertainty, these projections may affect the accuracy of TRY-based allocation results, especially for later years in the forecast horizon. Second, the feedback multiplier (0.25) for the dynamic TOPSIS allocation process was chosen by expert judgment rather than empirical optimization. Different values of the multiplier could affect the rate of change of priority classes across provinces, so testing the sensitivity of this parameter is another important avenue for future research.

Another limitation is that the proposed framework has not been compared to real budget allocation decisions made by the Ministry of Health. Therefore, further evaluation of its practical performance under real planning conditions is still needed.

\section{Conclusion}

\rev{This study developed an integrated ML framework for province-level
healthcare resource planning in Turkey for the 2025--2030 period, addressing
three interrelated problems: infrastructure forecasting, need-based provincial
ranking, and national budget projection. Walk-forward cross-validation across 19
models showed that time-series methods, particularly HoltWinters, ARIMA, and
Prophet, consistently achieved the highest accuracy for infrastructure targets,
while a weighted ensemble strategy was adopted to mitigate the flat-forecast
behaviour of ARIMA on sparse series. TOPSIS-based rankings identified small
northern provinces such as Sinop, Artvin, and Kastamonu as facing the greatest
healthcare need in 2025, driven by high elderly population ratios and low
physician density, with Shannon entropy assigning the strongest discriminating
power to the 65+ age ratio. The dynamic feedback mechanism prevented budget
monopolisation, shifting priority to larger underserved provinces such as
Balıkesir, Aydın, and Muğla from 2026 onward as early recipients reduced their
relative need scores. Budget forecasts were produced in USD to avoid distortions
from Turkish lira depreciation, using an ensemble of ExtraTrees, ARIMA, CAGR,
and GB models that achieved 88--90\% accuracy.}

\rev{The study makes three contributions. First, it provides a province-level
infrastructure forecast covering all 81 Turkish provinces across six indicators,
generated under a validation scheme that preserves temporal order. Second, it
introduces an entropy-weighted TOPSIS ranking that is recalculated annually in
response to allocations already made, so that priority circulates rather than
concentrating in a fixed set of provinces. Third, it links a national budget
forecast to province-level allocations, connecting macroeconomic projection to
sub-national distribution within a single framework.}

\rev{Several directions remain open. The feedback multiplier was set by expert
judgement and would benefit from validation-based optimisation. Long-term
exchange rate uncertainty could be represented through stochastic scenarios
rather than point projections. Comparing the framework's allocations against
actual Ministry of Health decisions would establish its practical value, and the
same three-module structure could be applied to other countries with
comparable provincial health statistics. The framework is designed to be
re-executed annually, so projections can be replaced by realised values as they
become available.}

\section*{Article Information}
\vspace{0.2cm}

\textbf{Author's Contributions:}
Bedirhan Köksoy: Software development, writing of the Method and Results sections. Ali Paşaoğlu: Supervision, writing of the Introduction and Conclusion sections. Şevval Duman: Software development, writing of the Method and Results sections. Sara Pournaji İran: Manuscript organization, English language editing, and writing the Literature Review section.\vspace{0.2cm}

\textbf{Conflict of Interest:}
This article contains no conflicts of interest.\vspace{0.2cm}

\textbf{Ethical Approval:}
The data used in this study do not contain any personal information and consist of official and approved data published by TURKSTAT. Therefore, the study does not violate ethical principles.
\vspace{0.2cm}

\textbf{Availability:}
The source code and related outputs are \href{https://github.com/koksoybedirhan/Forecasting-Provincial-Healthcare-Infrastructure}{available}.\vspace{0.2cm}

\textbf{Artificial Intelligence Statement:}
\rev{No artificial intelligence tools were used during the writing phase of this study.}
\vspace{0.2cm}

\textbf{Plagiarism Statement:}
This article has been scanned by iThenticate\texttrademark.
\vspace{0.2cm}

\begin{thebibliography}{99}

%%% Introduction

\bibitem{intro4} \rev{B. B. Alkan and A. Şahin, "Measuring inequalities in the distribution of health workers by bi-plot approach: The case of Turkey," \emph{J. Econ. Behav. Stud.}, vol. 2, no. 2, pp. 57--66, 2011, doi: 10.22610/jebs.v2i2.222.}

\bibitem{intro5} \rev{P. Abay, A. S. Yücel, M. Korkmaz, and N. Özdayı, "Health expenditure dynamics in Türkiye for the years 1999--2024 and their implications for the health system," \emph{Balıkesir J. Health Sci.}, vol. 15, no. 1, pp. 23--36, 2026.}

\bibitem{intro3} \rev{M. Isik, F. Isik, and M. Kiyak, "Analyzing financial structure of Turkish healthcare system in comparison with US, German, British, French and Cuban healthcare systems," \emph{J. Econ. Finance Account.}, vol. 2, no. 4, pp. 501--518, 2015.}

\bibitem{intro6} \rev{A. B. Şimşek, "Analysis of the change in the distribution of health human resources at the provincial level in Turkey in 2014--2019," \emph{Int. J. Health Manag. Strategies Res.}, vol. 8, no. 2, pp. 171--184, 2022.}

\bibitem{intro7} \rev{S. Söyler and D. Çavmak, "An evaluation of access to healthcare services in Türkiye," \emph{Aurum J. Health Sci.}, vol. 5, no. 2, pp. 69--76, 2023.}

\bibitem{intro8} \rev{D. T. Koç and E. Eren, "Workforce forecasting with machine learning for healthcare management," \emph{J. Health Manag.}, vol. 28, no. 1, pp. 121--135, 2026, doi: 10.1177/09720634251396339.}

\bibitem{intro9} \rev{M. Aslan and E. Toros, "Machine learning in optimising nursing care delivery models: An empirical analysis of hospital wards," \emph{J. Eval. Clin. Pract.}, vol. 31, no. 1, Feb. 2025, Art. no. e70001, doi: 10.1111/jep.70001.}

\bibitem{intro10} \rev{B. Meydan, "Evaluation of the performance of TR33 region provinces within the scope of sustainable development goals using the TOPSIS method" (in Turkish), \emph{J. Econ. Adm. Sci. Afyon Kocatepe Univ.}, vol. 27, Special Issue, pp. 227--249, 2025.}

\bibitem{lit3} M. Zhou, "Evaluation of health resource allocation efficiency and equity based on the entropy weight and TOPSIS methods," \emph{BMJ Open}, vol. 14, no. 3, 2024, Art. no. e082721, doi: 10.1136/bmjopen-2023-082721.

\bibitem{lit6} Y. Liu, H. Niu, X. Tian, M. Zhang, J. Feng, Y. Qian, et al., "Research on equity of medical resource allocation in Yangtze River Economic Belt under healthy China strategy," \emph{Front. Public Health}, vol. 11, 2023, Art. no. 1175276, doi: 10.3389/fpubh.2023.1175276.

\bibitem{lit4} J. Huang, Y. Wu, H. Pan, H. Yuan, X. Liu, H. Wang, et al., "Healthcare quality evaluation of tertiary public hospitals in ethnic border regions under China's performance assessment system: Based on the entropy weight TOPSIS method and RSR fuzzy set," \emph{Front. Public Health}, vol. 13, 2025, Art. no. 1668065, doi: 10.3389/fpubh.2025.1668065.

\bibitem{lit5} X. Dai, Y. Jiang, Y. Li, X. Wang, R. Wang, and Y. Zhang, "Evaluation of community basic public health service effect in a city in Inner Mongolia Autonomous Region: Based on entropy weight TOPSIS method and RSR fuzzy set," \emph{Arch. Public Health}, vol. 81, no. 1, 2023, Art. no. 149, doi: 10.1186/s13690-023-01151-x.

\bibitem{lit2} Y. Liu, L. Gong, H. Niu, F. Jiang, S. Du, and Y. Jiang, "Health system efficiency and equity in ASEAN: An empirical investigation," \emph{Cost Eff. Resour. Alloc.}, vol. 22, 2024, Art. no. 86, doi: 10.1186/s12962-024-00588-3.

\bibitem{lit1} G. Yeşilaydın, "Performance assessment of public hospitals with the entropy-weighted TOPSIS method: The case of Turkey," \emph{Hospital Topics}, vol. 104, no. 1, 2026, doi: 10.1080/00185868.2024.2404694.

\bibitem{lit7} C. Song, L. Fang, M. Xie, Z. Tang, Y. Zhang, and F. Tian, "Revealing spatiotemporal inequalities, hotspots, and determinants in healthcare resource distribution: Insights from hospital beds panel data in 2308 Chinese counties," \emph{BMC Public Health}, vol. 24, no. 1, 2024, Art. no. 423, doi: 10.1186/s12889-024-17950-y.

\bibitem{intro11} \rev{Z. Ceylan and A. Atalan, "Estimation of healthcare expenditure per capita of Turkey using artificial intelligence techniques with genetic algorithm-based feature selection," \emph{J. Forecast.}, vol. 40, no. 2, pp. 279--290, 2021, doi: 10.1002/for.2747.}

\bibitem{intro12} \rev{D. Güleryüz, "Predicting health spending in Turkey using the GPR, SVR, and DT models," \emph{Acta Infologica}, vol. 5, no. 1, pp. 155--166, 2021, doi: 10.26650/acin.882660.}

\bibitem{intro13} \rev{B. Köksoy, A. Paşaoğlu, and A. Akbel, "Forecasting regional healthcare workforce demand in Turkey using machine learning algorithms," in \emph{Proc. Int. Conf. Inf. Sci. Commun. Technol. (ICISCT)}, 2025, pp. 104--109.}

%%% Method

\bibitem{method1} \rev{D. Terzi Coban and O. Uslu, "Institutional performance evaluation at TUIK: A model proposal based on the balanced scorecard," \emph{Adm. Soc.}, vol. 57, no. 5, pp. 567--608, 2025, doi: 10.1177/00953997251328919.}

\bibitem{method2} \rev{S. García, J. Luengo, and F. Herrera, \emph{Data Preprocessing in Data Mining}, vol. 72. Cham, Switzerland: Springer, 2015, pp. 59--139, doi: 10.1007/978-3-319-10247-4.}

\bibitem{method3} \rev{Turkish Statistical Institute. \emph{Central Dissemination System (MEDAS)}. Accessed: Nov. 23, 2025. [Online]. Available: \url{https://biruni.tuik.gov.tr/medas}}

\bibitem{method4} \rev{T. Mustafy and M. T. U. Rahman, "Excel," in \emph{Statistics and Data Analysis for Engineers and Scientists}. Singapore: Springer Nature, 2024, pp. 81--134.}

\bibitem{method5} \rev{R. Mitchell, \emph{Web Scraping with Python: Collecting More Data from the Modern Web}, 2nd ed. Sebastopol, CA, USA: O'Reilly Media, 2018.}

\bibitem{methodnew1} \rev{M. A. Lones, "How to avoid machine learning pitfalls: A guide for academic researchers," 2021, \emph{arXiv:2108.02497}.}

\bibitem{methodnew2} \rev{M. Sakkari and M. Zaied, "A convolutional deep self-organizing map feature extraction for machine learning," \emph{Multimedia Tools Appl.}, vol. 79, no. 27, pp. 19451--19470, 2020, doi: 10.1007/s11042-020-08822-9.}

\bibitem{methodnew3} \rev{T. T. Wong and P. Y. Yeh, "Reliable accuracy estimates from k-fold cross validation," \emph{IEEE Trans. Knowl. Data Eng.}, vol. 32, no. 8, pp. 1586--1594, Aug. 2020, doi: 10.1109/TKDE.2019.2912815.}

\bibitem{methodnew4} \rev{I. K. Nti, O. Nyarko-Boateng, and J. Aning, "Performance of machine learning algorithms with different K values in K-fold cross-validation," \emph{Int. J. Inf. Technol. Comput. Sci.}, vol. 13, no. 6, pp. 61--71, 2021, doi: 10.5815/ijitcs.2021.06.05.}

\bibitem{methodnew5} \rev{D. M. Belete and M. D. Huchaiah, "Grid search in hyperparameter optimization of machine learning models for prediction of HIV/AIDS test results," \emph{Int. J. Comput. Appl.}, vol. 44, no. 9, pp. 875--886, 2022, doi: 10.1080/1206212X.2021.1974663.}

\bibitem{methodnew6} \rev{M. Feurer and F. Hutter, "Hyperparameter optimization," in \emph{Automated Machine Learning: Methods, Systems, Challenges}. Cham, Switzerland: Springer, 2019, pp. 3--33, doi: 10.1007/978-3-030-05318-5\_1.}

\bibitem{methodnew7} \rev{M. M. Ahsan, M. P. Mahmud, P. K. Saha, K. D. Gupta, and Z. Siddique, "Effect of data scaling methods on machine learning algorithms and model performance," \emph{Technologies}, vol. 9, no. 3, 2021, Art. no. 52, doi: 10.3390/technologies9030052.}

\bibitem{method6} \rev{Z. H. Zhou, \emph{Machine Learning}. Singapore: Springer Nature, 2021, doi: 10.1007/978-981-15-1967-3.}

\bibitem{method7} \rev{G. James, D. Witten, T. Hastie, R. Tibshirani, and J. Taylor, "Linear regression," in \emph{An Introduction to Statistical Learning: With Applications in Python}. Cham, Switzerland: Springer, 2023, pp. 69--134.}

\bibitem{method8} \rev{N. Herawati, A. Sutrisno, N. Nusyirwan, and M. Misgiyati, "The performance of ridge regression, LASSO, and elastic-net in controlling multicollinearity: A simulation and application," \emph{J. Mod. Appl. Stat. Methods}, vol. 23, no. 2, pp. 1--13, 2024.}

\bibitem{method10} \rev{L. Grinsztajn, E. Oyallon, and G. Varoquaux, "Why do tree-based models still outperform deep learning on typical tabular data?," in \emph{Adv. Neural Inf. Process. Syst.}, vol. 35, 2022, pp. 507--520.}

\bibitem{method11} \rev{U. J. Nzenwata, E. Edwin, E. A. Chukwu, D. Osilaja, J. O. Hinmikaiye, and C. Enyinnah, "Extra trees model for heart disease prediction," \emph{J. Data Anal. Inf. Process.}, vol. 13, no. 2, pp. 125--139, 2025.}

\bibitem{method12} \rev{J. T. Hancock and T. M. Khoshgoftaar, "CatBoost for big data: An interdisciplinary review," \emph{J. Big Data}, vol. 7, no. 1, 2020, Art. no. 94, doi: 10.1186/s40537-020-00369-8.}

\bibitem{method13} \rev{S. Sivaramakrishnan, T. F. Fernandez, R. G. Babukarthik, and S. Premalatha, "Forecasting time series data using ARIMA and Facebook Prophet models," in \emph{Big Data Management in Sensing}. Gistrup, Denmark: River Publishers, 2022, pp. 47--59.}

\bibitem{method14} \rev{A. J. U. Montiel, Á. A. Z. Morales, L. W. C. Cari, and J. R. M. Vergara, "Evolution of the SMEs loan portfolio in Ecuador: A Holt-Winters approach and the Extreme Learning Machine network" (in Spanish), \emph{Encuentros: Rev. Cienc. Humanas Teoría Soc. Pensam. Crít.}, no. 24, pp. 417--436, 2025.}

\bibitem{method15} \rev{E. Bauer and R. Kohavi, "An empirical comparison of voting classification algorithms: Bagging, boosting, and variants," \emph{Mach. Learn.}, vol. 36, no. 1, pp. 105--139, 1999, doi: 10.1023/A:1007515423169.}

\bibitem{method16} \rev{M. Shahhosseini, G. Hu, and H. Pham, "Optimizing ensemble weights and hyperparameters of machine learning models for regression problems," \emph{Mach. Learn. Appl.}, vol. 7, 2022, Art. no. 100251, doi: 10.1016/j.mlwa.2022.100251.}

\bibitem{method17} \rev{W. Zhou, S. Wang, M. Cucuringu, Z. Zhang, X. Li, J. Guo, et al., "DeltaLag: Learning dynamic lead-lag patterns in financial markets," in \emph{Proc. 6th ACM Int. Conf. AI Finance}, 2025, pp. 422--430.}

\bibitem{method18} \rev{I. Elbatal, M. Sarwar, F. Jamal, M. Daniyal, Z. Hussain, and A. B. Ghorbal, "Modelling on gross domestic product annual growth rate data by using time series, machine learning, and probability models," \emph{J. Radiat. Res. Appl. Sci.}, vol. 18, no. 2, 2025, Art. no. 101481, doi: 10.1016/j.jrras.2025.101481.}

\bibitem{methodnew8} \rev{T. G. Dietterich, "Ensemble methods in machine learning," in \emph{Proc. Int. Workshop Multiple Classifier Syst.}, Berlin, Germany: Springer, 2000, pp. 1--15, doi: 10.1007/3-540-45014-9\_1.}

\bibitem{methodnew9} \rev{Z. H. Zhou, \emph{Ensemble Methods: Foundations and Algorithms}, 2nd ed. Boca Raton, FL, USA: Chapman and Hall/CRC, 2025.}

\bibitem{methodnew10} \rev{D. Mishra, S. M. Tripathi, A. Chaurasia, and P. K. Chaurasia, "A review on ensemble learning methods: Machine learning approach," \emph{Int. J. Res. Publ. Rev.}, vol. 6, no. 2, pp. 3795--3803, 2025.}

\bibitem{method19} \rev{D. Chicco, M. J. Warrens, and G. Jurman, "The coefficient of determination R-squared is more informative than SMAPE, MAE, MAPE, MSE and RMSE in regression analysis evaluation," \emph{PeerJ Comput. Sci.}, vol. 7, 2021, Art. no. e623, doi: 10.7717/peerj-cs.623.}

\bibitem{method20} \rev{K. M. Sujon, R. Hassan, K. Choi, and M. A. Samad, "Accuracy, precision, recall, F1-score, or MCC? Empirical evidence from advanced statistics, ML, and XAI for evaluating business predictive models," \emph{J. Big Data}, vol. 12, no. 1, 2025, Art. no. 268, doi: 10.1186/s40537-025-01313-4.}

\bibitem{result1} \rev{R. J. Hyndman and G. Athanasopoulos, \emph{Forecasting: Principles and Practice}, 2nd ed. Melbourne, VIC, Australia: OTexts, 2018. Accessed: Nov. 23, 2025. [Online]. Available: \url{https://otexts.com/fpp2/}}

\bibitem{result2} \rev{A. C. Türkmen, T. Januschowski, Y. Wang, and A. T. Cemgil, "Forecasting intermittent and sparse time series: A unified probabilistic framework via deep renewal processes," \emph{PLOS ONE}, vol. 16, no. 11, 2021, Art. no. e0259764, doi: 10.1371/journal.pone.0259764.}

\bibitem{result3} \rev{J. Papathanasiou and N. Ploskas, "TOPSIS," in \emph{Multiple Criteria Decision Aid: Methods, Examples and Python Implementations} (Springer Optimization and Its Applications, vol. 136). Cham, Switzerland: Springer, 2018, pp. 1--30, doi: 10.1007/978-3-319-91648-4\_1.}

\bibitem{result4} \rev{P. Miç and Z. F. Antmen, "A healthcare facility location selection problem with fuzzy TOPSIS method for a regional hospital," \emph{Eur. J. Sci. Technol.}, no. 16, pp. 750--757, 2019, doi: 10.31590/ejosat.584217.}

\bibitem{result5} \rev{P. A. Bromiley, N. A. Thacker, and E. Bouhova-Thacker, "Shannon entropy, Renyi entropy, and information," \emph{Stat. Inf. Ser.}, vol. 9, pp. 2--8, 2004.}

\bibitem{result6} \rev{S. Vajapeyam, "Understanding Shannon's entropy metric for information," 2014, \emph{arXiv:1405.2061}.}

\bibitem{result7} \rev{S. A. Sepúlveda-Fontaine and J. M. Amigó, "Applications of entropy in data analysis and machine learning: A review," \emph{Entropy}, vol. 26, no. 12, 2024, Art. no. 1126, doi: 10.3390/e26121126.}

\end{thebibliography}
\end{document}